% BEGIN INCLUDED SGA7I_English_current_buildable_source_20260823/components/33_expose_VI_lemma_3_11.tex
% R3 explicit revisions VI-A-06; see qa/FIDELITY-R3-CHANGES.json.
\underline{引理 \SGARefTarget{sga7.sga7i.evi.n3-11}{3.11}.} 存在同构
\[
 D^2(A|R,M)\simeq
 \operatorname{Coker}\bigl(
 H^1(K_I,M)\longrightarrow
 \operatorname{Hom}_A(H_1(K_I),M)\bigr).
\]

\underline{证明.} 有
\begin{align*}
D^2(A|R,M)
 &=H^1\bigl(A|R,
   \underline H^1_{A|R}(\operatorname{Der}_R(-,M))\bigr)\\
 &=\ker\left(
   D^1(P_A^{(1)}|R,M)
   \mathrel{\substack{\longrightarrow\\[-.6ex]\longrightarrow\\[-.6ex]\longrightarrow}}
   D^1(P_A^{(2)}|R,M)\right),
\end{align*}
因为
\[
 [\underline H^1_{A|R}(\operatorname{Der}_R(-,M))](P_A^{(0)})
 =[\underline H^1_{A|R}(\operatorname{Der}_R(-,M))](P)=0.
\]

设
\[
 0\longrightarrow K\longrightarrow F\xrightarrow{\gamma}I
 \longrightarrow0
\]
为正合列. 记 $P[F]$ 为 $F$ 在 $P$ 上生成的对称代数 (symmetric algebra). 则 $P[F]$ 和
$P[F]\otimes_PP[F]$ 都是 $R$ 上的多项式代数, 且有满射
\begin{align*}
P[F]&\longrightarrow P_A^{(1)}=P\times_AP,
&X&\longmapsto(0,\gamma(X)),\\
P[F]\otimes_PP[F]&\longrightarrow
P_A^{(2)}=P\times_AP\times_AP,
&X\otimes1&\longmapsto(0,\gamma(X),0),\\[-.3ex]
&&1\otimes X&\longmapsto(0,0,\gamma(X)),
\end{align*}
其中 $X\in F$. 令
\[
 J=\ker(P[F]\longrightarrow P_A^{(1)}),
 \qquad
 J'=\ker(P[F]\otimes_PP[F]\longrightarrow P_A^{(2)}).
\]
若 $F_\gamma$ 表示由 $\{X-\gamma(X)\mid X\in F\}$ 生成的理想, 则
\begin{align*}
J&=(F)\cap(F_\gamma)=(K)+(F\cdot F_\gamma),\\
J'&=(F\otimes1+1\otimes F)
 \cap(F_\gamma\otimes1+1\otimes F)
 \cap(F\otimes1+1\otimes F_\gamma)\\
&=(K\otimes1+1\otimes K)
 +(F\cdot F_\gamma\otimes1+1\otimes F\cdot F_\gamma)
 +(F\otimes F).
\end{align*}

%%%%%%%%%%  scan idx 96  |  folio 85  %%%%%%%%%%

三个映射 $\pi_i:P_A^{(2)}\longrightarrow P_A^{(1)}$ 均有提升
$\widehat\pi_i:P[F]\otimes_PP[F]\longrightarrow P[F]$, 对 $X\in F$, 其定义如下:
\[
\begin{array}{c@{\qquad}c@{\qquad}c}
\widehat\pi_0:
\begin{cases}
 X\otimes1\longmapsto\gamma(X)-X,\\
 1\otimes X\longmapsto X,
\end{cases}
&
\widehat\pi_1:
\begin{cases}
 X\otimes1\longmapsto0,\\
 1\otimes X\longmapsto X,
\end{cases}
&
\widehat\pi_2:
\begin{cases}
 X\otimes1\longmapsto X,\\
 1\otimes X\longmapsto0.
\end{cases}
\end{array}
\]
直接计算可见, $\widehat\pi_0-\widehat\pi_1+\widehat\pi_2$ 在下列子模上为零:
\[
 K\otimes1+1\otimes K+1\otimes F\cdot F_\gamma
\]
并将 $F\cdot F_\gamma\otimes1+F\otimes F$ 满射到 $F\cdot F_\gamma$. 对每个
\[
 D\in\operatorname{Der}_R(P[F]\otimes_PP[F],M)
\]
均有 $D(F\cdot F_\gamma\otimes1+F\otimes F)=0$, 因为 $I\cdot M=0$. 因此,
\begin{align*}
D^2(A|R,M)
 &=\ker\left(
   D^1(P_A^{(1)}|R,M)
   \mathrel{\substack{\longrightarrow\\[-.6ex]\longrightarrow\\[-.6ex]\longrightarrow}}
   D^1(P_A^{(2)}|R,M)\right)\\
 &=\{[f]\in D^1(P_A^{(1)}|R,M)
       \mid f(F\cdot F_\gamma)=0\}\\
 &=\operatorname{Coker}\left(
   \operatorname{Der}_R(P[F],M)\longrightarrow
   \operatorname{Hom}_{P[F]}
   \left(\frac{(F)\cap(F_\gamma)}{F\cdot F_\gamma},M\right)
   \right)\\
 &=\operatorname{Coker}\left(
   \operatorname{Der}_P(P[F],M)\longrightarrow
   \operatorname{Hom}_{P[F]}
   \left(\frac{(F)\cap(F_\gamma)}{(F)(F_\gamma)},M\right)
   \right).
\end{align*}

下列序列
\[
 0\longrightarrow(F)\longrightarrow P[F]
 \xrightarrow{P[0]}P\longrightarrow0,
 \qquad
 0\longrightarrow(F_\gamma)\longrightarrow P[F]
 \xrightarrow{P[\gamma]}P_\gamma\longrightarrow0
\]
是正合的. 因而
\[
 \frac{(F)\cap(F_\gamma)}{(F)(F_\gamma)}
 =\operatorname{Tor}_1^{P[F]}(P,P_\gamma)
 =H_1(K_I),
\]
其中 $K_I$ 是与 $F\xrightarrow{\gamma}I\subset P$ 相伴的 Koszul 复形. 因此得到
\begin{align*}
D^2(A|R,M)
 &=\operatorname{Coker}\bigl(
   \operatorname{Hom}_P(F,M)\longrightarrow
   \operatorname{Hom}_P(H_1(K_I),M)\bigr)\\
 &=\operatorname{Coker}\bigl(
   H^1(K_I,M)\longrightarrow
   \operatorname{Hom}_A(H_1(K_I),M)\bigr).
\end{align*}
\hfill$\square$

现在可以把已经建立的结果归纳起来, 陈述下文所需的基本事实.

% END INCLUDED SGA7I_English_current_buildable_source_20260823/components/33_expose_VI_lemma_3_11.tex

% BEGIN INCLUDED SGA7I_English_current_buildable_source_20260823/components/34_expose_VI_theorem_3_12_statement.tex
%%%%%%%%%%  scan idx 97  |  folio 86  %%%%%%%%%%

\underline{定理 \SGARefTarget{sga7.sga7i.evi.n3-12}{3.12}.}

\begin{enumerate}[label=(\arabic*)]
\item 若
\[
 0\longrightarrow M'\longrightarrow M\longrightarrow M''
 \longrightarrow0
\]
为 $A$-模正合列, 则序列
\begin{multline*}
0\longrightarrow D^0(A|R,M')\longrightarrow D^0(A|R,M)
\longrightarrow D^0(A|R,M'')\longrightarrow D^1(A|R,M')\\
\longrightarrow D^1(A|R,M)\longrightarrow D^1(A|R,M'')
\longrightarrow\cdots\longrightarrow D^n(A|R,M')\\
\longrightarrow D^n(A|R,M)\longrightarrow D^n(A|R,M'')
\longrightarrow\cdots
\end{multline*}
是正合的.

\item 设 $S\longrightarrow R\longrightarrow A$ 为环同态. 则对每个 $A$-模 $M$, 序列
\begin{multline*}
0\longrightarrow D^0(A|R,M)\longrightarrow D^0(A|S,M)
\longrightarrow D^0(R|S,M)\longrightarrow D^1(A|R,M)\\
\longrightarrow D^1(A|S,M)\longrightarrow D^1(R|S,M)
\longrightarrow\cdots\longrightarrow D^n(A|R,M)\\
\longrightarrow D^n(A|S,M)\longrightarrow D^n(R|S,M)
\longrightarrow\cdots
\end{multline*}
是正合的.

\item 设 $A$ 和 $S$ 为 $R$-代数, 且对每个 $i>0$ 都有 $\operatorname{Tor}_i^R(A,S)=0$. 则对每个 $A\otimes_RS$-模 $M$, 典范映射
\[
 D^i(A\otimes_RS|S,M)\longrightarrow D^i(A|R,M)
\]
对每个 $i$ 都是同构. 特别地, 若 $A$ 在 $R$ 上平坦, 则
\[
 D^i(A\otimes_RS|S,M)\xrightarrow{\ \sim\ }D^i(A|R,M)
\]
对每个 $i$、每个 $R$-代数 $S$ 以及每个 $A\otimes_RS$-模 $M$ 都成立.

\item 设 $S$ 是 $A$ 的乘法子集. 对每个 $S^{-1}A$-模 $N$, 典范映射
\[
 D^i(S^{-1}A|R,N)\longrightarrow D^i(A|R,N)
\]
对每个 $i$ 都是同构. 从而对每个 $A$-模 $M$, 有典范同构
\[
 D^i(S^{-1}A|R,S^{-1}M)
 \xrightarrow{\ \sim\ }S^{-1}D^i(A|R,M)
\]
对每个 $i$ 成立.

\item 设 $A$ 和 $B$ 为 $R$-代数, 且当 $0<i\leq n$ 时 $\operatorname{Tor}_i^R(A,B)=0$. 则对每个 $A\otimes_RB$-模 $M$, 典范映射
\[
 D^i(A\otimes_RB|R,M)
 \longrightarrow D^i(A|R,M)\oplus D^i(B|R,M)
\]
当 $i\leq n$ 时是同构. 特别地, 假设对每个 $i>0$ 均有 $\operatorname{Tor}_i^R(A,A)=0$, 且对角态射
\[
 \Delta:\operatorname{Spec}(A)\longrightarrow
 \operatorname{Spec}(A)\times_{\operatorname{Spec}(R)}
 \operatorname{Spec}(A)
\]
为开浸入——例如, 假设 $A$ 在 $R$ 上\emph{平展}. 则
\[
 D^i(A|R,M)=0
\]
对每个 $i>0$ 及每个 $A$-模 $M$ 成立.

%%%%%%%%%%  scan idx 98  |  folio 87  %%%%%%%%%%

\item 设 $R$ 为 诺特环, $A$ 为本质有限型 $R$-代数, $E$ 为平坦 $A$-模. 则典范映射
\[
 E\otimes_AD^i(A|R,M)
 \xrightarrow{\ \sim\ }
 D^i(A|R,E\otimes_AM)
\]
对每个 $i\geq0$ 及每个 $A$-模 $M$ 都是同构.

\item 若 $R$ 为 诺特环, $A$ 为本质有限型 $R$-代数, $M$ 为有限型 $A$-模, 则对每个 $i$, $D^i(A|R,M)$ 均为有限型 $A$-模.

\item 设 $P$ 为 $R$ 上的多项式代数, 且
\[
 0\longrightarrow I\longrightarrow P\longrightarrow A
 \longrightarrow0
\]
为正合列. 则
\[
 D^0(A|R,M)=\operatorname{Der}_R(A,M),
\]
\begin{align*}
D^1(A|R,M)
 &=\operatorname{Coker}\bigl(
   \operatorname{Der}_R(P,M)\longrightarrow
   \operatorname{Hom}_A(I/I^2,M)\bigr)\\
 &=\left\{
   \begin{array}{c}
   \text{下列交换代数扩张的同构类:}\
   R\text{-代数扩张, 商为 }A\text{, 核为 }M
   \end{array}\right\},\\
D^2(A|R,M)
 &=\operatorname{Coker}\bigl(
   H^1(K_I,M)\longrightarrow
   \operatorname{Hom}_A(H_1(K_I),M)\bigr).
\end{align*}
\end{enumerate}

% END INCLUDED SGA7I_English_current_buildable_source_20260823/components/34_expose_VI_theorem_3_12_statement.tex

% BEGIN INCLUDED SGA7I_English_current_buildable_source_20260823/components/35_expose_VI_theorem_3_12_proof.tex
\underline{证明.}

\begin{enumerate}[label=(\arabic*)]
\item 若
\[
 0\longrightarrow M'\longrightarrow M\longrightarrow M''
 \longrightarrow0
\]
为 $A$-模正合列, 则对 $R$ 上的每个多项式代数 $P$, 序列
\[
 0\longrightarrow\operatorname{Der}_R(P,M')
 \longrightarrow\operatorname{Der}_R(P,M)
 \longrightarrow\operatorname{Der}_R(P,M'')
 \longrightarrow0
\]
是正合的. 因而, 由 \SGARefLink{sga7.sga7i.evi.n3-3}{3.3}(c),
\[
 0\longrightarrow\operatorname{Der}_R(-,M')
 \longrightarrow\operatorname{Der}_R(-,M)
 \longrightarrow\operatorname{Der}_R(-,M'')
 \longrightarrow0
\]
为层的正合列.

\item 函子 $\operatorname{Der}_S(-,M)$ 是 $\underline C_{A|S}$ 上的加性层, 且
\[
 \operatorname{Der}_R(B,M)
 =\ker\bigl(\operatorname{Der}_S(B,M)
 \longrightarrow\operatorname{Der}_S(R,M)\bigr)
\]
对每个 $B\in\operatorname{Ob}(\underline C_{A|R})$ 都成立. 因此, 断言由下述结果得出: \SGARefLink{sga7.sga7i.evi.n3-7}{3.7}(c).

\item 设
\[
 S:\underline C_{A|R}\longrightarrow
 \underline C_{A\otimes_RS|S},
 \qquad X\longmapsto X\otimes_RS.
\]
于是 $\widetilde S\operatorname{Der}_S(-,M)
=\operatorname{Der}_R(-,M)$, 因而断言由下述结果得出: \SGARefLink{sga7.sga7i.evi.n3-7}{3.7}(a).

%%%%%%%%%%  scan idx 99  |  folio 88  %%%%%%%%%%

\item 此结论来自 \SGARefLink{sga7.sga7i.evi.n3-9}{3.9}. 事实上, $F=\operatorname{Der}_R(-,N)$ 是 $\underline C_{S^{-1}A|R}$ 上的加性层, 而 $F_A\in\operatorname{Ob}
(\widetilde{\underline C}_{S^{-1}A|A})$ 也是加性的, 因为
\begin{align*}
 F_A(X)
 &=\ker\bigl(F(X)\longrightarrow F(A)\bigr)\\
 &=\ker\bigl(\operatorname{Der}_R(X,N)
   \longrightarrow\operatorname{Der}_R(A,N)\bigr)\\
 &=\operatorname{Der}_A(X,N).
\end{align*}
故 $F_A=\operatorname{Der}_A(-,N)$.

\item 第一个断言来自 \SGARefLink{sga7.sga7i.evi.n3-8}{3.8}(a). 现假设对每个 $i>0$ 均有 $\operatorname{Tor}_i^R(A,A)=0$, 且
\[
 \Delta:\operatorname{Spec}(A)\longrightarrow
 \operatorname{Spec}(A)\times_{\operatorname{Spec}(R)}
 \operatorname{Spec}(A)
\]
为开浸入. 设 $\varphi:A\otimes_RA\longrightarrow A$ 为乘法映射. 对 $A$ 的每个极大理想 $\mathfrak m$, 诱导映射
\[
 (A\otimes_RA)_{\varphi^{-1}(\mathfrak m)}
 \longrightarrow A_{\mathfrak m}
\]
是同构. 因而由第一个断言及 \SGARefLink{sga7.sga7i.evi.n3-12}{3.12}(4) 得到, 对角映射
\[
 \Delta_{\mathfrak m}:D^i(A|R,M)_{\mathfrak m}
 \longrightarrow D^i(A|R,M)_{\mathfrak m}
 \oplus D^i(A|R,M)_{\mathfrak m}
\]
对 $A$ 的每个极大理想 $\mathfrak m$ 都是同构. 从而,
\[
 \Delta:D^i(A|R,M)\longrightarrow
 D^i(A|R,M)\oplus D^i(A|R,M)
\]
是同构, 故对每个 $i>0$ 都有 $D^i(A|R,M)=0$.

\item 由 \SGARefLink{sga7.sga7i.evi.n3-12}{3.12}(4), 可假设 $A$ 为有限型 $R$-代数. 考虑层态射
\[
 E\otimes_A\operatorname{Der}_R(-,M)
 \longrightarrow
 \operatorname{Der}_R(-,E\otimes_AM)
\]
定义在 $\widetilde{\underline C}_{A|R}$ 上. 由于 $E$ 在 $A$ 上平坦, 映射
\[
 E\otimes_A\operatorname{Der}_R(B,M)
 \xrightarrow{\ \sim\ }
 \operatorname{Der}_R(B,E\otimes_AM)
\]
只要 $B\in\operatorname{Ob}(\underline C_{A|R})$ 在 $R$ 上为有限型便是同构. 由 \SGARefLink{sga7.sga7i.evi.n3-3}{3.3}(b) 可知,
\[
 D^i\bigl(A|R,E\otimes_A\operatorname{Der}_R(-,M)\bigr)
 \longrightarrow
 D^i\bigl(A|R,\operatorname{Der}_R(-,E\otimes_AM)\bigr)
\]
是同构. 另一方面, $E\otimes_A-$ 是正合的, 因此
\[
 E\otimes_AD^i(A|R,M)
 =D^i\bigl(A|R,E\otimes_A\operatorname{Der}_R(-,M)\bigr).
\]
断言得证.

\item 此结论来自 \SGARefLink{sga7.sga7i.evi.n3-3}{3.3}(a).

%%%%%%%%%%  scan idx 100  |  folio 89  %%%%%%%%%%

\item 由 \SGARefLink{sga7.sga7i.evi.n3-10.occ-2}{3.10} 和 \SGARefLink{sga7.sga7i.evi.n3-11}{3.11}, 只需证明
\[
 \operatorname{Coker}\bigl(
 \operatorname{Der}_R(P,M)\longrightarrow
 \operatorname{Hom}_A(I/I^2,M)\bigr)
 =\operatorname{Exalcomm}(A|R,M),
\]
其中 $\operatorname{Exalcomm}(A|R,M)$ 表示以 $M$ 为核的 $A$ 的 $R$-代数扩张的同构类集合.
设
\[
 0\longrightarrow M\longrightarrow A'\longrightarrow A
 \longrightarrow0
\]
为这样的扩张. 由于 $P$ 为 $R$ 上的多项式代数, 得到交换图
\[
\begin{tikzcd}[column sep=large]
0 \arrow[r] & I \arrow[r] \arrow[d]
& P \arrow[r] \arrow[d]
& A \arrow[r] \arrow[d,equal] & 0 \\
0 \arrow[r] & M \arrow[r]
& A' \arrow[r]
& A \arrow[r] & 0.
\end{tikzcd}
\]
映射 $I\longrightarrow M$ 在相差 $\operatorname{Der}_R(P,M)$ 的作用的意义下由扩张 $A'$ 决定. 因此得到映射
\[
 \Phi:\operatorname{Exalcomm}(A|R,M)
 \longrightarrow
 \operatorname{Coker}\bigl(
 \operatorname{Der}_R(P,M)\longrightarrow
 \operatorname{Hom}_A(I/I^2,M)\bigr).
\]
反之, 一个 $P$-线性映射 $I\longrightarrow M$ 通过推出给出扩张, 嵌入下图:
\[
\begin{tikzcd}[column sep=large]
0 \arrow[r] & I \arrow[r] \arrow[d]
& P \arrow[r] \arrow[d]
& A \arrow[r] \arrow[d,equal] & 0 \\
0 \arrow[r] & M \arrow[r]
& A' \arrow[r]
& A \arrow[r] & 0.
\end{tikzcd}
\]
两个相差 $\operatorname{Der}_R(P,M)$ 中一个元素的 $P$-线性映射 $I\rightrightarrows M$ 给出同构的扩张. 因而有映射
\[
 \Psi:\operatorname{Coker}\bigl(
 \operatorname{Der}_R(P,M)\longrightarrow
 \operatorname{Hom}_A(I/I^2,M)\bigr)
 \longrightarrow\operatorname{Exalcomm}(A|R,M).
\]
容易验证 $\Phi$ 与 $\Psi$ 互为逆映射.
\end{enumerate}
\hfill$\square$

% END INCLUDED SGA7I_English_current_buildable_source_20260823/components/35_expose_VI_theorem_3_12_proof.tex

% BEGIN INCLUDED SGA7I_English_current_buildable_source_20260823/components/36_expose_VI_corollary_3_13.tex
%%%%%%%%%%  scan idx 100  |  folio 89  %%%%%%%%%%

\underline{推论 \SGARefTarget{sga7.sga7i.evi.n3-13}{3.13}.}

\begin{enumerate}[label=(\alph*)]
\item $R$-代数 $A$ 在 $R$ 上形式光滑, 当且仅当对每个 $A$-模 $M$ 都有 $D^1(A|R,M)=0$.

\item 设 $R$ 为 诺特环, $A$ 为有限型 $R$-代数. 则 $A$ 在 $R$ 上为相对完全交, 当且仅当对每个 $A$-模 $M$ 都有 $D^2(A|R,M)=0$.

\item 设 $R$ 为 诺特环, $A$ 为有限型 $R$-代数. 对 $A$ 的每个乘法子集 $S$, 存在典范同构
\[
 D^i(S^{-1}A|R,S^{-1}M)
 \xrightarrow{\ \sim\ }S^{-1}D^i(A|R,M)
 \qquad(i\geq0).
\]

\item 设 $R$ 为 诺特环, $A$ 为本质有限型局部 $R$-代数. 记 $\widehat A$ 为 $A$ 的 $\mathfrak m$-进完备化, 其中 $\mathfrak m$ 为 $A$ 的极大理想. 则对每个有限型 $A$-模 $E$, 存在典范同构
\[
 D^1(\widehat A|R,\widehat A\otimes_AE)
 \xrightarrow{\ \sim\ }
 \widehat A\otimes_AD^1(A|R,E).
\]
\end{enumerate}

\underline{证明.}

(a) 立即由 $D^1(A|R,M)=\operatorname{Exalcomm}(A|R,M)$ 得出.

对于 (b), 设 $P$ 为 $R$ 上有限个变元的多项式代数, 且
\[
 0\longrightarrow I\longrightarrow P\longrightarrow A
 \longrightarrow0
\]
为正合列. 局部化后可假设 $I$ 由一个 $P$-正则序列生成. 于是 $H_1(K_I)=0$, 故
\[
 D^2(A|R,M)
 =\operatorname{Coker}\bigl(
 H^1(K_I,M)\longrightarrow
 \operatorname{Hom}_A(H_1(K_I),M)\bigr)=0.
\]
反之, 假设对每个 $A$-模 $M$ 都有 $D^2(A|R,M)=0$. 取正合列
\[
 0\longrightarrow K\longrightarrow F\longrightarrow I
 \longrightarrow0,
\]
其中 $F$ 为有限型自由 $P$-模. 若
\[
 K_0=\operatorname{Im}\bigl(\textstyle\bigwedge^2F
 \longrightarrow\bigwedge^1F\bigr),
\]
则 $H_1(K_I)=K/K_0$. 由于 $D^2(A|R,H_1(K_I))=0$, 映射
\[
 H^1(K_I,H_1(K_I))
 \longrightarrow
 \operatorname{Hom}_A(H_1(K_I),H_1(K_I))
\]
为满射. 因此, 正合列
\[
 0\longrightarrow K/K_0\longrightarrow F/IF
 \longrightarrow I/I^2\longrightarrow0
\]
作为 $A$-模序列是可裂的. 特别地, $I/I^2$ 为投射模. 再对 $A$ 局部化后, 可假设 $I/I^2$ 自由, 且 $F/IF\longrightarrow I/I^2$ 为同构. 于是由上述可裂正合列得到 $H_1(K_I)=K/K_0=0$. 故 $I$ 由一个 $P$-正则序列生成, 从而 $A$ 为 $R$ 上的相对完全交.

(c) 立即由下述结果得出: \SGARefLink{sga7.sga7i.evi.n3-12}{3.12}(4) 和 (6).

%%%%%%%%%%  scan idx 101  |  folio 90  %%%%%%%%%%

对于 (d), \SGARefLink{sga7.sga7i.evi.n3-12}{3.12}(6) 首先给出典范同构
\[
 D^i(A|R,\widehat A\otimes_AE)
 \xrightarrow{\ \sim\ }
 \widehat A\otimes_AD^i(A|R,E)
\]
对每个 $i$ 成立. 因而只需证明
\[
 D^1(\widehat A|R,\widehat A\otimes_AE)
 \longrightarrow D^1(A|R,\widehat A\otimes_AE)
\]
为同构. 我们使用 $D^1$ 关于代数扩张的描述.

设
\[
 0\longrightarrow\widehat A\otimes_AE\longrightarrow A'
 \xrightarrow{\ \varphi\ }\widehat A\longrightarrow0
\]
为一个限制到 $A$ 后可裂的 $R$-代数扩张. 由于 $\widehat A\otimes_AE$ 为有限型 $\widehat A$-模, $A'$ 为局部 $R$-代数, 且关于其 $\mathfrak m'$-进拓扑完备, 其中 $\mathfrak m'$ 为 $A'$ 的极大理想. 若 $\psi:A\longrightarrow A'$ 是提升典范映射 $A\longrightarrow\widehat A$ 的 $R$-代数映射, 则 $\psi$ 唯一延拓为映射 $\widehat\psi:\widehat A\longrightarrow A'$, 满足 $\varphi\circ\widehat\psi=\operatorname{id}_{\widehat A}$. 这就证明了
\[
 D^1(\widehat A|R,\widehat A\otimes_AE)
 \longrightarrow D^1(A|R,\widehat A\otimes_AE)
\]
为单射.

%%%%%%%%%%  scan idx 102  |  folio 91  %%%%%%%%%%

现在设
\[
 0\longrightarrow\widehat A\otimes_AE\longrightarrow B
 \xrightarrow{\ \varphi\ }A\longrightarrow0
\]
为 $R$-代数扩张, 并令 $\widehat E=\widehat A\otimes_AE$. 首先断言, $B$ 上的 $\mathfrak m_B$-进拓扑在 $\widehat E$ 上诱导的拓扑, 与其作为 $\widehat A$-模的拓扑一致. 事实上, 由于 $A$ 在 $R$ 上为本质有限型, 存在子代数 $B_0\subset B$, 使得 $B_0$ 为本质有限型局部 $R$-代数, 且 $\varphi(B_0)=A$. 令 $E_0=\widehat E\cap B_0$, 分别以 $\mathfrak m_B$、$\mathfrak m_0$、$\mathfrak m$ 表示 $B$、$B_0$、$A$ 的极大理想. 则
\[
 \mathfrak m_B=\mathfrak m_0+\widehat E,
 \qquad
 \mathfrak m_B^{n+1}=\mathfrak m_0^{n+1}
   +\mathfrak m^n\widehat E
\]
对每个 $n$ 成立. 由于 $B_0$ 为 诺特环, Artin--Rees 定理给出整数 $r>0$, 使得对每个 $n\geq r$, 有
\begin{align*}
 \widehat E\cap\mathfrak m_B^{n+1}
 &=\mathfrak m^n\widehat E
   +\widehat E\cap\mathfrak m_0^{n+1}\\
 &=\mathfrak m^n\widehat E
   +E_0\cap\mathfrak m_0^{n+1}\\
 &=\mathfrak m^n\widehat E
   +\mathfrak m^{n+1-r}(E_0\cap\mathfrak m_0^r)
 \subset\mathfrak m^{n+1-r}\widehat E.
\end{align*}
这就证明了关于这两种拓扑的断言. 特别地, $\widehat E$ 关于 $B$ 的 $\mathfrak m_B$-进拓扑所诱导的拓扑完备, 因而完备化给出行正合的交换图
\[
\begin{tikzcd}[column sep=large]
0 \arrow[r] & \widehat E \arrow[r] \arrow[d,equal]
& B \arrow[r] \arrow[d]
& A \arrow[r] \arrow[d] & 0 \\
0 \arrow[r] & \widehat E \arrow[r]
& \widehat B \arrow[r]
& \widehat A \arrow[r] & 0.
\end{tikzcd}
\]
这里 $\widehat B$ 表示 $B$ 的 $\mathfrak m_B$-进完备化. 因此
\[
 D^1(\widehat A|R,\widehat E)
 \longrightarrow D^1(A|R,\widehat E)
\]
为满射, 证明完毕.
\hfill$\square$

% END INCLUDED SGA7I_English_current_buildable_source_20260823/components/36_expose_VI_corollary_3_13.tex

% BEGIN INCLUDED SGA7I_English_current_buildable_source_20260823/components/37_expose_VI_3_14_through_3_16.tex
%%%%%%%%%%  scan idx 103  |  folio 92  %%%%%%%%%%

\underline{推论 \SGARefTarget{sga7.sga7i.evi.n3-14}{3.14}.}
设 $R$ 为 诺特环, $A$ 为有限型 $R$-代数, 且除一个闭点 $x\in\operatorname{Spec}(A)$ 外处处光滑. 则有同构
\[
 D^1(A|R,A)\simeq D^1(A_x|R,A_x)
 \simeq D^1(\widehat A_x|R,\widehat A_x).
\]

\underline{证明.}
由于 $A$ 在 $x$ 以外在 $R$ 上光滑, $D^1(A|R,A)$ 是支集包含在 $\{x\}$ 中的有限型 $A$-模. 因此
\[
 D^1(A|R,A)=D^1(A|R,A)_x
 =\widehat{D^1(A|R,A)_x}.
\]
现由下述结果即得断言: \SGARefLink{sga7.sga7i.evi.n3-13}{3.13}.
\hfill$\square$

\underline{注 \SGARefTarget{sga7.sga7i.evi.n3-15}{3.15}.}
设 $R$ 为拓扑环, $A$ 为拓扑 $R$-代数. 对每个 $A$-模 $E$, 定义
\[
 D^i_{\mathrm{top}}(A|R,E)
 =\varinjlim_\alpha
 D^i(A_\alpha|R_\alpha,A_\alpha\otimes_AE),
\]
其中 $A_\alpha=A/J_\alpha$, $R_\alpha=R/I_\alpha$, 且 $J_\alpha$ 和 $I_\alpha$ 分别为 $A$ 和 $R$ 中的开理想. 因而 $A$ 作为拓扑 $R$-代数形式光滑, 当且仅当对每个 $A$-模 $E$ 都有 $D^1_{\mathrm{top}}(A|R,E)=0$.

现在设 $X$ 为环 $R$ 上的概形, $\underline E$ 为拟凝聚 $\mathcal O_X$-模. 由仿射开子集 $U\subset X$ 上的预层定义 $X$ 上的层
\[
 \underline D^i_{X|R}(\underline E)
 =\underline D^i(X|R,\underline E)
\]
该预层给定为
\[
 \bigl[\underline D^i_{X|R}(\underline E)\bigr](U)
 =D^i\bigl(\Gamma(U,\mathcal O_X)|R,
            \Gamma(U,\underline E)\bigr).
\]
若 $R$ 为 诺特环, $X$ 在 $R$ 上为有限型, 则局部化性质 \SGARefLink{sga7.sga7i.evi.n3-13}{3.13}(c) 表明, 这些预层典范地构成拟凝聚 $\mathcal O_X$-模.

在不致混淆时, 简记
\[
 \underline D^i_{X|R}
 =\underline D^i_{X|R}(\mathcal O_X).
\]
\SGARefLink{sga7.sga7i.evi.n3-13}{3.13} 的所有断言, 对拟凝聚层 $\underline D^i_{X|R}(\underline E)$ 都有相应的表述. 这些表述留给读者, 下面只列出与我们的目的最相关的性质.

\underline{推论 \SGARefTarget{sga7.sga7i.evi.n3-15.occ-2}{3.15}.}
设 $R$ 为 诺特环, $X$ 为有限型 $R$-概形.

%%%%%%%%%%  scan idx 104  |  folio 93  %%%%%%%%%%

\begin{enumerate}[label=(\alph*)]
\item 对每个凝聚 $\mathcal O_X$-模 $\underline E$ 及每个 $i$, $\underline D^i_{X|R}(\underline E)$ 仍为凝聚模. 此外,
\[
 \operatorname{Supp}\underline D^i_{X|R}(\underline E)
 \subset X^{\mathrm{Sing}}
 \qquad(i>0),
\]
其中
\[
 X^{\mathrm{Sing}}
 =\{x\in X\mid X\text{ 在 }R\text{ 上不光滑于 }x\}.
\]

\item 若 $X$ 在 $R$ 上平坦, 则对每个 $R$-代数 $S$, 存在典范同构
\[
 \underline D^i_{X\otimes_RS\,|\,S}(\underline E)
 \xrightarrow{\ \sim\ }
 \underline D^i_{X|R}(\underline E)
\]
对每个 $i$ 成立.

\item 假设 $X$ 为 $R$ 上的相对完全交. 对每个满环同态 $S\longrightarrow R$, 有正合列
\[
0\longrightarrow\underline D^1_{X|R}(\underline E)
\longrightarrow\underline D^1_{X|S}(\underline E)
\longrightarrow
\underline{\operatorname{Hom}}_{\mathcal O_X}
 \bigl((I/I^2)\otimes_R\mathcal O_X,\underline E\bigr)
\longrightarrow0,
\]
其中 $I=\ker(S\longrightarrow R)$. 若进一步假设 $R$ 为局部环, 剩余域为 $k$, 且 $X$ 在 $R$ 上平坦, 则
\[
0\longrightarrow
 \underline D^1_{X_0|k}(\mathcal O_{X_0})
\longrightarrow
 \underline D^1_{X|S}(\mathcal O_{X_0})
\longrightarrow
 D^1(R|S,k)\otimes_k\mathcal O_{X_0}
\longrightarrow0,
\]
其中 $X_0=X\otimes_Rk$.
\end{enumerate}

\underline{证明.}
(a) 由 \SGARefLink{sga7.sga7i.evi.n3-12}{3.12} 中的有限性和低次描述得到, 而 (b) 由 \SGARefLink{sga7.sga7i.evi.n3-12}{3.12}(3) 得到.

对于 (c), 由于 $X$ 为 $R$ 上的相对完全交, 对每个 $i\geq2$ 都有 $\underline D^i_{X|R}(\underline E)=0$. 因而传递性序列 \SGARefLink{sga7.sga7i.evi.n3-12}{3.12}(2) 给出正合列
\begin{multline*}
0\longrightarrow\underline D^0_{X|R}(\underline E)
\longrightarrow\underline D^0_{X|S}(\underline E)
\longrightarrow\underline D^0_{R|S}(\underline E)
\longrightarrow\underline D^1_{X|R}(\underline E)\\
\longrightarrow\underline D^1_{X|S}(\underline E)
\longrightarrow\underline D^1_{R|S}(\underline E)
\longrightarrow0,
\end{multline*}
其中 $X$ 上的拟凝聚层 $\underline D^i_{R|S}(\underline E)$ 在每个仿射开子集 $U\subset X$ 上定义为
\[
 \bigl[\underline D^i_{R|S}(\underline E)\bigr](U)
 =D^i(R|S,\underline E(U)).
\]
由于 $S\longrightarrow R$ 为满射, $\underline D^0_{R|S}(\underline E)=0$; 又由 \SGARefLink{sga7.sga7i.evi.n3-12}{3.12}(8) 中的低次描述得
\[
 \underline D^1_{R|S}(\underline E)
 =\underline{\operatorname{Hom}}_{\mathcal O_X}
  \bigl((I/I^2)\otimes_R\mathcal O_X,\underline E\bigr).
\]
这就证明了第一个正合列.

%%%%%%%%%%  scan idx 105  |  folio 94  %%%%%%%%%%

若 $X$ 在 $R$ 上平坦, 则对每个拟凝聚 $\mathcal O_{X_0}$-模 $\underline E$, 有典范同构
\[
 \underline D^i_{X_0|k}(\underline E)
 \xrightarrow{\ \sim\ }
 \underline D^i_{X|R}(\underline E)
\]
对每个 $i$ 成立. 特别地,
\[
 \underline D^1_{X_0|k}(\mathcal O_{X_0})
 \xrightarrow{\ \sim\ }
 \underline D^1_{X|R}(\mathcal O_{X_0}),
\]
第二个正合列由此得出.
\hfill$\square$

\underline{注 \SGARefTarget{sga7.sga7i.evi.n3-16}{3.16}.}
上面定义的拟凝聚层 $\underline D^i_{X|R}(\underline E)$ 应当是联系局部理论与整体理论的某个谱序列中的项 $E_2^{0,i}$. L.~Illusie 已在 \emph{Lecture Notes in Mathematics} 239 中建立了这样的理论. 他的方法使用“同伦代数”, 推广了 M.~André 和 D.~Quillen 的方法, 而非使用“上同调代数”. 看来, 可以在概形范畴上赋予适当的拓扑, 推广仿射范畴上的拓扑, 并利用其上的 Grothendieck 上同调, 将这里的方法整体化.

% END INCLUDED SGA7I_English_current_buildable_source_20260823/components/37_expose_VI_3_14_through_3_16.tex

% BEGIN INCLUDED SGA7I_English_current_buildable_source_20260823/components/38_expose_VI_section_4_opening_through_4_1.tex
%%%%%%%%%%  scan idx 105  |  folio 94  %%%%%%%%%%

\section*{\SGARefTarget{sga7.sga7i.evi.n4}{4}. 形变的形式模}

在一些自然的余纤维范畴中, 对象的同构类一般不可 pro 表, 但仍有形式遍有对象; 形变的余纤维范畴就是其中之一. 现在回到第 1 节的情形与记号. 因此, 在本节中始终固定一个剩余域为 $k$ 的完备 Noether 局部环 $\Lambda$, 并以 $\underline C_\Lambda$ 表示剩余域为 $k$ 的 Artin 局部 $\Lambda$-代数所成的范畴.

设 $X_0$ 为固定的 $k$ 上的概形. 所谓 $X_0$ 的一个 (无穷小) 形变 (deformation), 是指一对 $(R,X)$, 其中 $R\in\operatorname{Ob}(\underline C_\Lambda)$, $X$ 为平坦 $R$-概形, 再加上一个交换图
\[
\begin{tikzcd}[column sep=huge,row sep=huge]
X_0 \arrow[r,"i_X"] \arrow[d] & X \arrow[d,"\text{平坦}"] \\
\operatorname{Spec}(k) \arrow[r] & \operatorname{Spec}(R)
\end{tikzcd}
\]
使得诱导态射
\[
 X_0\longrightarrow
 X\times_{\operatorname{Spec}(R)}\operatorname{Spec}(k)
\]
是同构. 由于 $R$ 为 Artin 局部环, $i_X:X_0\longrightarrow X$ 必为闭浸入, 且在底拓扑空间上为同胚.

%%%%%%%%%%  scan idx 106  |  folio 95  %%%%%%%%%%

若 $(R,X)$ 和 $(R',X')$ 分别是 $X_0$ 在 $R$ 和 $R'$ 上的形变, 则箭头
\[
 (R',X')\longrightarrow(R,X)
\]
由 $\underline C_\Lambda$ 中的映射 $\varphi:R'\longrightarrow R$ 以及概形态射 $\Phi:X\longrightarrow X'$ 组成, 它们嵌入交换图
\[
\begin{tikzcd}[column sep=huge,row sep=large]
X \arrow[rr,"\Phi"] \arrow[ddd] && X' \arrow[ddd] \\
& X_0 \arrow[ul,"i_X"'] \arrow[ur,"i_{X'}"] \arrow[d] & \\
& \operatorname{Spec}(k) \arrow[dl] \arrow[dr] & \\
\operatorname{Spec}(R)
 \arrow[rr,"\operatorname{Spec}(\varphi)"']
&& \operatorname{Spec}(R').
\end{tikzcd}
\]
这就定义了 $X_0$ 的无穷小形变范畴 $\underline M_{X_0}$. 它带有函子
\[
 p:\underline M_{X_0}\longrightarrow\underline C_\Lambda,
 \qquad (R,X)\longmapsto R.
\]

设 $(R,X)\in\operatorname{Ob}(\underline M_{X_0})$, $\varphi:R\longrightarrow S$ 为 $\underline C_\Lambda$ 中的映射. 有交换图
\[
\begin{tikzcd}[column sep=huge,row sep=huge]
X\times_{\operatorname{Spec}(R)}\operatorname{Spec}(S)
 \arrow[r,"p_1"] \arrow[d,"\text{平坦}"']
& X \arrow[d,"\text{平坦}"] \\
\operatorname{Spec}(S)
 \arrow[r,"\operatorname{Spec}(\varphi)"']
& \operatorname{Spec}(R),
\end{tikzcd}
\]
而 $\varphi$-箭头
\[
 (\varphi,p_1):(R,X)\longrightarrow
 \bigl(S,X\times_{\operatorname{Spec}(R)}\operatorname{Spec}(S)\bigr)
\]
在 $\underline M_{X_0}$ 中为余笛卡尔箭头. 因此 $\underline M_{X_0}$ 是 $\underline C_\Lambda$ 上的余纤维范畴. 下面列出若干事实.

%%%%%%%%%%  scan idx 107  |  folio 96  %%%%%%%%%%

\medskip
\noindent\textbf{\SGARefTarget{sga7.sga7i.evi.n4-1.item-a}{\SGARefAliasTarget{sga7.sga7i.evi.n4-1}4.1}(a).}
设
\[
 (1_R,\Phi):(R,X_1)\longrightarrow(R,X_2)
\]
为 $\underline M_{X_0}$ 中的箭头. 于是 $\Phi:X_2\longrightarrow X_1$ 为 $R$-态射, 其约化
\[
 X_2\otimes_Rk\longrightarrow X_1\otimes_Rk
\]
为同构. 由于 $R$ 为 Artin 环且各 $X_i$ 在 $R$ 上平坦, $\Phi$ 为同构. 因此 $\underline M_{X_0}$ 是 $\underline C_\Lambda$ 上的余纤维群胚.

\medskip
\noindent\textbf{\SGARefTarget{sga7.sga7i.evi.n4-1.item-b}{4.1}(b).}
假设 $X_0$ 为仿射概形, 记作 $X_0=\operatorname{Spec}(A_0)$. 则对每个 $(R,X)\in\operatorname{Ob}(\underline M_{X_0})$, $X$ 也为仿射概形. 因此, 对每个 $R\in\operatorname{Ob}(\underline C_\Lambda)$, 取整体截面将 $\underline M_{X_0}(R)$ 等同于如下范畴. 一个对象是平坦 $R$-代数 $A$, 连同 $R$-代数映射 $j_A:A\longrightarrow A_0$, 后者诱导同构
\[
 A\otimes_Rk\xrightarrow{\ \sim\ }A_0.
\]
态射 $\Phi:A_1\longrightarrow A_2$ 是使下图交换的 $R$-代数映射:
\[
\begin{tikzcd}[column sep=huge,row sep=huge]
A_1 \arrow[r,"\Phi"] \arrow[d,"j_{A_1}"']
& A_2 \arrow[d,"j_{A_2}"] \\
A_0 \arrow[r,equal] & A_0
\end{tikzcd}
\]
此图交换.

\medskip
\noindent\textbf{\SGARefTarget{sga7.sga7i.evi.n4-1.item-c}{4.1}(c).}
考虑 $\widehat{\underline C}_\Lambda$ 上的余纤维群胚 $\widehat{\underline M}_{X_0}$. 按定义, 对每个 $R\in\operatorname{Ob}(\widehat{\underline C}_\Lambda)$, 范畴 $\widehat{\underline M}_{X_0}(R)$ 典范等价于如下形式概形所成的范畴: $\mathfrak X$ 在 $\operatorname{Spf}(R)$ 上为进性形式概形, 且
\[
 \mathfrak X_n
 =\mathfrak X\otimes_R R/\mathfrak m^{n+1}
\]
对每个 $n\geq0$ 都在 $R/\mathfrak m^{n+1}$ 上平坦. 若 $X_0$ 为 诺特概形, 这等价于要求 $\mathfrak X$ 在 $\operatorname{Spf}(R)$ 上平坦. 因而在这种情形下, $\widehat{\underline M}_{X_0}$ 在 $\widehat{\underline C}_\Lambda$ 上典范等价于形式概形意义下的 $X_0$ 在 $\widehat{\underline C}_\Lambda$ 的对象上的形变范畴.

\medskip
\noindent\textbf{\SGARefTarget{sga7.sga7i.evi.n4-1.item-d}{4.1}(d).}
设 $(R,X)\in\operatorname{Ob}(\underline M_{X_0})$, 于是有交换图
\[
\begin{tikzcd}[column sep=huge,row sep=huge]
X_0 \arrow[r,"i_X"] \arrow[d] & X \arrow[d,"\text{平坦}"] \\
\operatorname{Spec}(k) \arrow[r] & \operatorname{Spec}(R)
\end{tikzcd}
\]
其诱导映射 $X_0\longrightarrow X\times_{\operatorname{Spec}(R)}
\operatorname{Spec}(k)$ 为同构.

%%%%%%%%%%  scan idx 108  |  folio 97  %%%%%%%%%%

设 $U_0\subset X_0$ 为开子概形. 由于 $i_X:X_0\longrightarrow X$ 为闭浸入, 且在底空间上为同胚, 限制给出交换图
\[
\begin{tikzcd}[column sep=huge,row sep=huge]
U_0 \arrow[r,"i_X|_{U_0}"] \arrow[d]
& X|_{U_0} \arrow[d,"\text{平坦}"] \\
\operatorname{Spec}(k) \arrow[r] & \operatorname{Spec}(R).
\end{tikzcd}
\]
诱导映射
\[
 U_0\longrightarrow
 (X|_{U_0})\times_{\operatorname{Spec}(R)}\operatorname{Spec}(k)
 =(X\otimes_Rk)|_{U_0}
\]
是同构, 故 $(R,X|_{U_0})$ 是 $U_0$ 在 $R$ 上的形变. 此外, $\underline M_{X_0}$ 中的每个箭头 $(\varphi,\Phi):(R,X)\longrightarrow(R',X')$ 都限制为 $\underline M_{U_0}$ 中的箭头
\[
 (\varphi,\Phi|_{U_0}):(R,X|_{U_0})
 \longrightarrow(R',X'|_{U_0})
\]
因此有限制函子
\[
 \underline M_{X_0}\longrightarrow\underline M_{U_0}
\]
在 $\underline C_\Lambda$ 上.

同样可以在一点处作局部化与完备化. 设 $x_0\in X_0$. 若 $(R,X)\in\operatorname{Ob}(\underline M_{X_0})$, 则 $\operatorname{Spec}(\mathcal O_{X,x_0})$ 是 $\operatorname{Spec}(\mathcal O_{X_0,x_0})$ 在 $R$ 上的形变, 而 $\operatorname{Spec}(\widehat{\mathcal O}_{X,x_0})$ 是 $\operatorname{Spec}(\widehat{\mathcal O}_{X_0,x_0})$ 在 $R$ 上的形变; 帽号表示关于极大理想各次幂定义的拓扑的完备化. 因此有函子
\[
 \underline M_{X_0}
 \longrightarrow
 \underline M_{\operatorname{Spec}(\mathcal O_{X_0,x_0})}
 \longrightarrow
 \underline M_{\operatorname{Spec}(\widehat{\mathcal O}_{X_0,x_0})}.
\]

首先证明 $\underline M_{X_0}$ 是 $\underline C_\Lambda$ 上的齐性余纤维范畴.

% END INCLUDED SGA7I_English_current_buildable_source_20260823/components/38_expose_VI_section_4_opening_through_4_1.tex

% BEGIN INCLUDED SGA7I_English_current_buildable_source_20260823/components/39_expose_VI_lemma_4_2_and_corollary_4_3.tex
%%%%%%%%%%  scan idx 108  |  folio 97  %%%%%%%%%%

\medskip
\noindent\textbf{引理 \SGARefTarget{sga7.sga7i.evi.n4-2}{4.2}.}
考虑 $\underline C_\Lambda$ 中的图
\[
\begin{tikzcd}[column sep=huge,row sep=large]
& R' \arrow[d] \\
R \arrow[r,"\text{满射}"'] & \bar R
\end{tikzcd}
\]
设 $E$ 和 $E'$ 分别为 $R$ 和 $R'$ 上的平坦模. 则对每个 $R'$-同态
\[
 E'\longrightarrow \bar E=E\otimes_R\bar R,
\]
纤维积 $E\times_{\bar E}E'$ 作为 $R\times_{\bar R}R'$ 上的模是平坦的.

%%%%%%%%%%  scan idx 109  |  folio 98  %%%%%%%%%%

\medskip
\noindent\textit{证明.}
设
\[
 0\longrightarrow I\longrightarrow R\longrightarrow\bar R
 \longrightarrow0
\]
为正合列, 从而
\[
 0\longrightarrow I\longrightarrow R\times_{\bar R}R'
 \longrightarrow R'\longrightarrow0
\]
为正合列. 由于 $E$ 在 $R$ 上平坦,
\[
 0\longrightarrow I\otimes_RE\longrightarrow E
 \longrightarrow\bar E\longrightarrow0
\]
为正合列, 因而得到正合列
\begin{equation}\tag{*}
0\longrightarrow I\otimes_RE\longrightarrow
E\times_{\bar E}E'\longrightarrow E'\longrightarrow0.
\end{equation}
置 $S=R\times_{\bar R}R'$. 由于 $I$ 幂零, $E\times_{\bar E}E'$ 作为 $S$-模平坦, 当且仅当 $(E\times_{\bar E}E')\otimes_SR'$ 作为 $R'$-模平坦, 且
\[
 \operatorname{Tor}_1^S(E\times_{\bar E}E',R')=0.
\]
而
\[
 I\otimes_RE=I\otimes_S(E\times_{\bar E}E'),
 \qquad
 E'=R'\otimes_S(E\times_{\bar E}E'),
\]
因而正合列 $(*)$ 可改写为
\[
0\longrightarrow I\otimes_S(E\times_{\bar E}E')
\longrightarrow E\times_{\bar E}E'
\longrightarrow R'\otimes_S(E\times_{\bar E}E')
\longrightarrow0.
\]
故
\[
 \operatorname{Tor}_1^S(R',E\times_{\bar E}E')=0,
 \qquad
 R'\otimes_S(E\times_{\bar E}E')=E',
\]
而 $E'$ 在 $R'$ 上平坦. 可知 $E\times_{\bar E}E'$ 在 $S$ 上平坦. \hfill$\square$

\medskip
\noindent\textbf{推论 \SGARefTarget{sga7.sga7i.evi.n4-3}{4.3}.}
给定 $\underline M_{X_0}$ 中的图
\[
\begin{tikzcd}[column sep=huge,row sep=large]
(R,X) \arrow[dr,"{(\varphi,\Phi)}"']
&& (R',X') \arrow[dl,"{(\varphi',\Phi')}"] \\
& (\bar R,\bar X) &
\end{tikzcd}
\]
若 $\varphi:R\longrightarrow\bar R$ 为满射, 则 (概形范畴中的) 融合和 $X\amalg_{\bar X}X'$ 是 $R\times_{\bar R}R'$ 上的平坦概形. 因而 $\underline M_{X_0}$ 是 $\underline C_\Lambda$ 上的齐性余纤维群胚.

\medskip
\noindent\textit{证明.}
对每个仿射开子集 $U_0\subset X_0$, 按定义有
\[
 (X\amalg_{\bar X}X')|_{U_0}
 \simeq
 \operatorname{Spec}\!\left(
   \Gamma(U_0,X)
   \times_{\Gamma(U_0,\bar X)}
   \Gamma(U_0,X')
 \right).
\]
因此, 若 $\varphi:R\longrightarrow\bar R$ 为满射, 则由引理 \SGARefLink{sga7.sga7i.evi.n4-2}{4.2}, $X\amalg_{\bar X}X'$ 在 $R\times_{\bar R}R'$ 上平坦. 此外, 态射
\[
 i_X\amalg_{i_{\bar X}}i_{X'}:
 X_0\longrightarrow X\amalg_{\bar X}X'
\]
诱导同构
\[
 X_0\longrightarrow (X\amalg_{\bar X}X')\otimes_Sk,
 \qquad S=R\times_{\bar R}R'.
\]
因此
\[
 \bigl(R\times_{\bar R}R',X\amalg_{\bar X}X'\bigr)
 \in\operatorname{Ob}(\underline M_{X_0}),
\]
且在 $\underline M_{X_0}$ 中有
\[
 \bigl(R\times_{\bar R}R',X\amalg_{\bar X}X'\bigr)
 =
 (R,X)\times_{(\bar R,\bar X)}(R',X').
\]
故 $\underline M_{X_0}$ 在 $\underline C_\Lambda$ 上为齐性的. \hfill$\square$

%%%%%%%%%%  scan idx 110  |  folio 99  %%%%%%%%%%

% END INCLUDED SGA7I_English_current_buildable_source_20260823/components/39_expose_VI_lemma_4_2_and_corollary_4_3.tex

% BEGIN INCLUDED SGA7I_English_current_buildable_source_20260823/components/40_expose_VI_lemma_4_4.tex
%%%%%%%%%%  scan idx 110  |  folio 99  %%%%%%%%%%

\medskip
\noindent\textbf{引理 \SGARefTarget{sga7.sga7i.evi.n4-4}{4.4}.}
对 $k$ 上的每个概形 $X_0$, 存在典范正合列
\[
0\longrightarrow H^1(X_0,\underline D^0_{X_0})
\longrightarrow[\underline M_{X_0}(k[\varepsilon])]
\longrightarrow H^0(X_0,\underline D^1_{X_0}),
\]
其中各 $\underline D^i_{X_0}$ 的定义见第 3 节.

\medskip
\noindent\textit{证明.}
首先假设 $X_0$ 为仿射概形, 令 $\Gamma(X_0,\mathcal O_{X_0})=A_0$. 如 \SGARefLink{sga7.sga7i.evi.n4-1.item-b}{4.1}(b) 所指出, 范畴 $\underline M_{X_0}(k[\varepsilon])$ 典范等价于一个范畴 $\underline E$, 其对象为平坦 $k[\varepsilon]$-代数 $A$, 连同诱导同构
\[
k\otimes_{k[\varepsilon]}A\xrightarrow{\ \sim\ }A_0.
\]
的 $k[\varepsilon]$-代数映射 $j_A:A\longrightarrow A_0$. 等价地, $\underline E$ 的一个对象是 $k$-代数 $A$, 连同正合列
\[
0\longrightarrow A_0\xrightarrow{\ \varepsilon\ }A
\xrightarrow{\ j_A\ }A_0\longrightarrow0
\]
其中 $\varepsilon(A_0)^2=0$ 在 $A$ 的理想意义下成立.

在第一种描述中, $\underline E$ 中的箭头 $A_1\longrightarrow A_2$ 是满足 $j_{A_1}=j_{A_2}\circ\Phi$ 的 $k[\varepsilon]$-代数映射 $\Phi:A_1\longrightarrow A_2$. 在第二种描述中, 它是给出如下交换图的 $k$-代数映射 $\Phi:A_1\longrightarrow A_2$:
\[
\begin{tikzcd}[column sep=large,row sep=large]
0 \arrow[r]
& A_0 \arrow[r,"\varepsilon"] \arrow[d,equal]
& A_1 \arrow[r,"j_{A_1}"] \arrow[d,"\Phi"]
& A_0 \arrow[r] \arrow[d,equal]
& 0 \\
0 \arrow[r]
& A_0 \arrow[r]
& A_2 \arrow[r,"j_{A_2}"]
& A_0 \arrow[r]
& 0.
\end{tikzcd}
\]
因此, 通过 $\Gamma(X_0,-)$, $[\underline M_{X_0}(k[\varepsilon])]$ 典范同构于以 $A_0$ 为核的 $A_0$ 的 $k$-代数扩张的同构类集合. 换言之, 有典范双射
\[
[\underline M_{X_0}(k[\varepsilon])]
\xrightarrow{\ \sim\ }H^0(X_0,\underline D^1_{X_0})
\]
它显然保持两侧的向量空间结构.

%%%%%%%%%%  scan idx 111  |  folio 100  %%%%%%%%%%

一般地, 对每个仿射开子集 $U_0\subset X_0$, 限制与仿射情形给出复合映射
\[
[\underline M_{X_0}(k[\varepsilon])]
\longrightarrow[\underline M_{U_0}(k[\varepsilon])]
\xrightarrow{\ \sim\ }H^0(U_0,\underline D^1_{X_0}).
\]
这些映射给出典范映射
\[
[\underline M_{X_0}(k[\varepsilon])]
\longrightarrow H^0(X_0,\underline D^1_{X_0}).
\]
它的核由 $X_0$ 在 $k[\varepsilon]$ 上的局部平凡形变的同构类组成, 即局部同构于 $X_0\otimes_k k[\varepsilon]$ 的形变的同构类. $X_0\otimes_k k[\varepsilon]$ 的在 $X_0$ 上诱导恒等的自同构所成的层为
\[
\underline D^0_{X_0}
=\underline{\operatorname{Der}}_k
  (\mathcal O_{X_0},\mathcal O_{X_0}).
\]
因此, 此类局部平凡形变的同构类典范地等同于 $H^1(X_0,\underline D^0_{X_0})$, 这就给出了所需的正合列. \hfill$\square$

% END INCLUDED SGA7I_English_current_buildable_source_20260823/components/40_expose_VI_lemma_4_4.tex

% BEGIN INCLUDED SGA7I_English_current_buildable_source_20260823/components/41_expose_VI_theorem_4_5.tex
% R3 explicit revisions VI-B-C03; see qa/FIDELITY-R3-CHANGES.json.
%%%%%%%%%%  scan idx 111  |  folio 100  %%%%%%%%%%

\medskip
\noindent\textbf{定理 \SGARefTarget{sga7.sga7i.evi.n4-5}{4.5}.}
设 $X_0$ 为 $k$ 上的有限型概形, 满足下列条件之一:

\smallskip
\noindent(i) $X_0$ 在 $k$ 上紧合; 或

\noindent(ii) $X_0$ 为仿射概形, 且
\[
X_0^{\mathrm{sing}}
=\{x\in X_0\mid X_0\text{ 在 }k\text{ 上不光滑于 }x\}
\]
为有限集.

\smallskip
\noindent\textbf{(a)}
范畴 $\underline M_{X_0}$ 在 $\underline C_\Lambda$ 上有形式遍有对象, 称为 $X_0$ 的形式遍有形变. 若 $(R,\mathcal X)$ 是 $X_0$ 的形式遍有形变, 则对每个 $(S,\mathcal Y)\in\operatorname{Ob}(\widehat{\underline M}_{X_0})$, 均存在 $\widehat{\underline M}_{X_0}$ 中的映射
\[
(\varphi,\Phi):(R,\mathcal X)\longrightarrow(S,\mathcal Y)
\]
使下图为笛卡尔图:
\[
\begin{tikzcd}[column sep=huge,row sep=large]
\mathcal Y \arrow[r,"\Phi"] \arrow[d]
& \mathcal X \arrow[d] \\
\operatorname{Spf}(S)
 \arrow[r,"\operatorname{Spf}(\varphi)"']
& \operatorname{Spf}(R).
\end{tikzcd}
\]
此外, $X_0$ 的形式遍有形变 $(R,\mathcal X)$ pro 表函子 $[\underline M_{X_0}]$, 当且仅当对 $\underline C_\Lambda$ 中的每个满射 $R'\longrightarrow R$ 以及 $X_0$ 在 $R'$ 上的每个形变 $(R',X')$, 同态
\[
\operatorname{Aut}_{R'}(X'\mid X_0)
\longrightarrow
\operatorname{Aut}_R(X'\otimes_{R'}R\mid X_0)
\]
为满射.

\smallskip
\noindent\textbf{(b)}
若 $X_0$ 在 $k$ 上紧合, 则 $\underline M_{X_0}$ 在 $\underline C_\Lambda$ 上光滑伪可表.

%%%%%%%%%%  scan idx 112  |  folio 101  %%%%%%%%%%

\medskip
\noindent\textit{证明.}
\textbf{(a)} 由推论 4.3, $\underline M_{X_0}$ 是 $\underline C_\Lambda$ 上的齐性余纤维群胚. 由 \SGARefLink{sga7.sga7i.evi.n1-11}{1.11} 以及 \SGARefLink{sga7.sga7i.evi.n2-9}{2.9}(b), 只需证明 $[\underline M_{X_0}(k[\varepsilon])]$ 在 $k$ 上为有限维.

考虑正合列
\begin{equation}\tag{*}
0\longrightarrow H^1(X_0,\underline D^0_{X_0})
\longrightarrow[\underline M_{X_0}(k[\varepsilon])]
\longrightarrow H^0(X_0,\underline D^1_{X_0}).
\end{equation}
若 $X_0$ 在 $k$ 上紧合, 则 $H^1(X_0,\underline D^0_{X_0})$ 和 $H^0(X_0,\underline D^1_{X_0})$ 均为有限维, 因为由 \SGARefLink{sga7.sga7i.evi.n3-15.occ-2}{3.15}(a), 对所有 $i\geq0$, $\underline D^i_{X_0}$ 都是凝聚层.

现在假设 $X_0$ 为仿射概形, 且 $X_0^{\mathrm{sing}}$ 为有限集. 层 $\underline D^1_{X_0}$ 为凝聚层, 且
\[
\operatorname{Supp}\underline D^1_{X_0}
\subset X_0^{\mathrm{sing}}
\]
这是 \SGARefLink{sga7.sga7i.evi.n3-15.occ-2}{3.15}(a) 的结论. 因而 $H^0(X_0,\underline D^1_{X_0})$ 在 $k$ 上为有限维, 而 $H^1(X_0,\underline D^0_{X_0})=0$, 因为 $X_0$ 仿射. 两种情形下, 正合列 $(*)$ 都表明 $[\underline M_{X_0}(k[\varepsilon])]$ 在 $k$ 上为有限维. (a) 的第二个断言来自 \SGARefLink{sga7.sga7i.evi.n2-7}{2.7}.

\smallskip
\noindent\textbf{(b)} 若 $X_0$ 紧合, 则
\[
\operatorname{Aut}(X_0\otimes_k k[\varepsilon]\mid X_0)
=H^0(X_0,\underline D^0_{X_0})
\]
为有限维, 断言由 \SGARefLink{sga7.sga7i.evi.n2-13}{2.13} 得出.
\hfill$\square$

% END INCLUDED SGA7I_English_current_buildable_source_20260823/components/41_expose_VI_theorem_4_5.tex

% BEGIN INCLUDED SGA7I_English_current_buildable_source_20260823/components/42_expose_VI_remarks_4_6_and_pair_deformations.tex
%%%%%%%%%%  scan idx 112  |  folio 101  %%%%%%%%%%

\medskip
\noindent\textbf{注 \SGARefTarget{sga7.sga7i.evi.n4-6.item-a}{\SGARefAliasTarget{sga7.sga7i.evi.n4-6}4.6}(a).}
设 $X_0$ 仿射. $X_0$ 存在形式遍有形变, 当且仅当 $\underline D^1_{X_0}$ 的支集为有限集. 然而这并不意味着 $X_0$ 的所有非光滑点都是孤立的: 存在 $k$ 上的仿射概形 $X_0$, 有非孤立的非光滑点, 但 $\operatorname{Supp}\underline D^1_{X_0}$ 为有限集.

\medskip
\noindent\textbf{注 \SGARefTarget{sga7.sga7i.evi.n4-6.item-b}{4.6}(b).}
今后, $X_0$ 的形变指 $\underline M_{X_0}$ 的对象. 在可能混淆时, 称这样的对象为 $X_0$ 的无穷小形变.

\medskip
\noindent\textbf{注 \SGARefTarget{sga7.sga7i.evi.n4-6.item-c}{4.6}(c).}
设 $X_0$ 为 $k$ 上的概形. 我们固定了剩余域为 $k$ 的完备 Noether 局部环 $\Lambda$, 并考虑了 $X_0$ 在 $\underline C_\Lambda$ 上的形变.

%%%%%%%%%%  scan idx 113  |  folio 102  %%%%%%%%%%

实际中通常取 $\Lambda=k$, 或在 $k$ 为正特征完美域时取 $\Lambda=W(k)$. 包含 $\underline C_k\subset\underline C_\Lambda$ 是全忠实的, 而 $\underline M_{X_0}|_{\underline C_k}$ 正是 $X_0$ 在 $\underline C_k$ 的对象上的形变范畴. 因此, 若 $(R,\mathfrak X)$ 是 $X_0$ 在 $\underline C_\Lambda$ 上的形式遍有形变, 则
\[
 (k\otimes_\Lambda R,\;k\otimes_\Lambda\mathfrak X)
\]
是 $X_0$ 在 $\underline C_k$ 上的形式遍有形变. 例如, 由此可知
\[
 \dim_{\underline C_\Lambda}\underline M_{X_0}
 =\dim(k\otimes_\Lambda R)
\]
仅依赖于 $X_0$, 与 $\Lambda$ 的选取无关.

前一定理有许多推广. 例如, 考虑一对 $Y_0\longrightarrow X_0$, 其中 $X_0$ 是 $k$ 上的有限型概形, $Y_0$ 为 $X_0$ 的固定闭子概形. 对 $R\in\operatorname{Ob}(\underline C_\Lambda)$, $(X_0,Y_0)$ 在 $R$ 上的形变是一个交换图
\[
\begin{tikzcd}[row sep=large,column sep=huge]
Y_0 \arrow[rrr] \arrow[dr,hook] \arrow[ddr]
&&& Y \arrow[dl,hook] \arrow[ddl,"\text{平坦}"] \\
& X_0 \arrow[r] \arrow[d]
& X \arrow[d,"\text{平坦}"'] & \\
& \operatorname{Spec}(k) \arrow[r]
& \operatorname{Spec}(R) &
\end{tikzcd}
\]
使得诱导映射
\[
(X_0,Y_0)\longrightarrow
\left(
 X\times_{\operatorname{Spec}(R)}\operatorname{Spec}(k),
 Y\times_{\operatorname{Spec}(R)}\operatorname{Spec}(k)
\right)
\]
为同构. 以显然的方式定义 $(X_0,Y_0)$ 的形变之间的态射
\[
(R';X',Y')\longrightarrow(R;X,Y)
\]
便得到 $\underline C_\Lambda$ 上的余纤维群胚 $\underline M_{X_0,Y_0}$. 这种构造经常用于形变模问题中, 以刚化自同构. 推论 \SGARefLink{sga7.sga7i.evi.n4-3}{4.3} 中基于引理 \SGARefLink{sga7.sga7i.evi.n4-2}{4.2} 的论证立即表明, $\underline M_{X_0,Y_0}$ 是 $\underline C_\Lambda$ 上的齐性余纤维群胚.

% END INCLUDED SGA7I_English_current_buildable_source_20260823/components/42_expose_VI_remarks_4_6_and_pair_deformations.tex

% BEGIN INCLUDED SGA7I_English_current_buildable_source_20260823/components/43_expose_VI_lemma_4_7_through_theorem_4_9.tex
%%%%%%%%%%  scan idx 113  |  folio 102  %%%%%%%%%%

\medskip
\noindent\textbf{引理 \SGARefTarget{sga7.sga7i.evi.n4-7}{4.7}.}
序列
\[
0\longrightarrow
H^0\!\left(
Y_0,
\underline{\operatorname{Hom}}_{\mathcal O_{Y_0}}
(\underline I/\underline I^2,\mathcal O_{Y_0})
\right)
\longrightarrow[\underline M_{X_0,Y_0}(k[\varepsilon])]
\longrightarrow[\underline M_{X_0}(k[\varepsilon])]
\]
为正合列, 其中 $\underline I$ 是在 $X_0$ 中定义 $Y_0$ 的理想层; 因而
\[
0\longrightarrow\underline I\longrightarrow\mathcal O_{X_0}
\longrightarrow\mathcal O_{Y_0}\longrightarrow0
\]
为正合列.

%%%%%%%%%%  scan idx 114  |  folio 103  %%%%%%%%%%

\medskip
\noindent\textit{证明.}
设 $\underline B$ 为 $\underline M_{X_0,Y_0}(k[\varepsilon])$ 的全子范畴, 其对象具有如下形式:
\[
\bigl(k[\varepsilon];
X_0\times_{\operatorname{Spec}(k)}\operatorname{Spec}(k[\varepsilon]),
Y\bigr).
\]
于是
\[
[\underline B]
=\ker\!\left(
[\underline M_{X_0,Y_0}(k[\varepsilon])]
\longrightarrow[\underline M_{X_0}(k[\varepsilon])]
\right).
\]
范畴 $\underline B$ 是 $Y_0$ 在 $X_0\times_{\operatorname{Spec}(k)}\operatorname{Spec}(k[\varepsilon])$ 内的 $k[\varepsilon]$ 上形变所成的群胚. 因而其对象是一个交换图
\[
\begin{tikzcd}[column sep=huge,row sep=large]
Y \arrow[rr,"i_Y"] \arrow[dr,"\text{平坦}"']
&& X_0\times_{\operatorname{Spec}(k)}
   \operatorname{Spec}(k[\varepsilon]) \arrow[dl] \\
& \operatorname{Spec}(k[\varepsilon]) &
\end{tikzcd}
\]
其特殊纤维为固定的闭浸入 $i_{Y_0}$. 箭头 $(Y_1,i_{Y_1})\longrightarrow(Y_2,i_{Y_2})$ 是 $k[\varepsilon]$-态射 $\Phi:Y_1\longrightarrow Y_2$ (必为同构), 其特殊纤维为 $\operatorname{id}_{Y_0}$, 且 $i_{Y_2}\circ\Phi=i_{Y_1}$.

等价地, $\underline B$ 的一个对象是一对 $(Y,j_Y)$, 其中 $Y$ 是 $Y_0$ 在 $k[\varepsilon]$ 上的形变, $j_Y:Y\longrightarrow X_0$ 是 $k$-态射, 使复合 $Y_0\longrightarrow Y\xrightarrow{j_Y}X_0$ 为 $i_{Y_0}$. 箭头 $(Y_1,j_{Y_1})\longrightarrow(Y_2,j_{Y_2})$ 是特殊纤维为 $\operatorname{id}_{Y_0}$ 且满足 $j_{Y_2}\circ\Phi=j_{Y_1}$ 的态射 $\Phi:Y_1\longrightarrow Y_2$.

将 $(Y,j_Y)$ 对应于图
\[
\begin{tikzcd}[column sep=large,row sep=large]
0 \arrow[r]
& \underline I \arrow[r] \arrow[d]
& \mathcal O_{X_0} \arrow[r] \arrow[d,"j_Y"]
& \mathcal O_{Y_0} \arrow[r] \arrow[d,equal]
& 0 \\
0 \arrow[r]
& \mathcal O_{Y_0} \arrow[r,"\varepsilon"]
& \mathcal O_Y \arrow[r]
& \mathcal O_{Y_0} \arrow[r]
& 0
\end{tikzcd}
\]
的对应给出同构
\[
\begin{aligned}
[\underline B]
&\xrightarrow{\ j_Y|_{\underline I}\ }
\operatorname{Hom}_{\mathcal O_{X_0}}
 (\underline I,\mathcal O_{Y_0}) \\
&=
\operatorname{Hom}_{\mathcal O_{Y_0}}
 (\underline I/\underline I^2,\mathcal O_{Y_0}) \\
&=
H^0\!\left(
Y_0,
\underline{\operatorname{Hom}}_{\mathcal O_{Y_0}}
(\underline I/\underline I^2,\mathcal O_{Y_0})
\right).
\end{aligned}
\]
正合性得证. \hfill$\square$

%%%%%%%%%%  scan idx 115  |  folio 104  %%%%%%%%%%

\medskip
\noindent\textbf{注 \SGARefTarget{sga7.sga7i.evi.n4-8}{4.8}.}
一个蕴含引理 \SGARefLink{sga7.sga7i.evi.n4-7}{4.7} 的更一般公式见 SGA 3, III, 4.4. 还可注意到如下事实. 对每个 $S$-概形 $X$ 及每个拟凝聚 $\mathcal O_X$-模 $E$, 都可以对所有 $i\geq0$ 定义拟凝聚 $\mathcal O_X$-模 $\underline D^i_{X|S}(E)$, 推广 $S$ 仿射时的情形. 与往常一样, 记
\[
\underline D^i_{X|S}(\mathcal O_X)=\underline D^i_{X|S}.
\]
于是, 对闭浸入 $Y_0\hookrightarrow X_0$, 有
\[
H^0\!\left(
Y_0,
\underline{\operatorname{Hom}}_{\mathcal O_{Y_0}}
(\underline I/\underline I^2,\mathcal O_{Y_0})
\right)
=H^0(Y_0,\underline D^1_{Y_0|X_0}),
\]
其中 $0\longrightarrow\underline I\longrightarrow\mathcal O_{X_0}
\longrightarrow\mathcal O_{Y_0}\longrightarrow0$ 为正合列.

\medskip
\noindent\textbf{定理 \SGARefTarget{sga7.sga7i.evi.n4-9}{4.9}.}
设 $X_0$ 为有限型 $k$-概形, $Y_0$ 为 $X_0$ 的闭子概形. 假设满足下列条件之一:

\smallskip
\noindent(i) $X_0$ 在 $k$ 上紧合; 或

\noindent(ii) $X_0$ 仿射, 且除有限多个点外在 $k$ 上光滑, 而 $Y_0$ 在 $k$ 上紧合.

\smallskip
\noindent\textbf{(a)}
范畴 $\underline M_{X_0,Y_0}$ 在 $\underline C_\Lambda$ 上有形式遍有形变. 它 pro 表函子 $[\underline M_{X_0,Y_0}]$, 当且仅当对 $\underline C_\Lambda$ 中的每个满射 $R'\longrightarrow R$ 以及每个形变 $(R';X',Y')$, 映射
\[
\operatorname{Aut}_{R'}((X',Y')\mid X_0)
\longrightarrow
\operatorname{Aut}_{R}\!\left(
(X'\otimes_{R'}R,Y'\otimes_{R'}R)\mid X_0
\right)
\]
为满射.

\smallskip
\noindent\textbf{(b)}
若此外
\[
H^0\!\left(
X_0,
\underline D^0_{Y_0|k}
\times_{\underline D^0_{X_0|k}(\mathcal O_{Y_0})}
\underline D^0_{X_0|k}
\right)
\]
为有限维, 则 $\underline M_{X_0,Y_0}$ 光滑伪可表.

\medskip
\noindent\textit{证明.}
由引理 \SGARefLink{sga7.sga7i.evi.n4-4}{4.4}、\SGARefLink{sga7.sga7i.evi.n4-7}{4.7} 以及 \SGARefLink{sga7.sga7i.evi.n1-11}{1.11}、\SGARefLink{sga7.sga7i.evi.n2-7}{2.7}、\SGARefLink{sga7.sga7i.evi.n2-13}{2.13} 立即得出. \hfill$\square$

% END INCLUDED SGA7I_English_current_buildable_source_20260823/components/43_expose_VI_lemma_4_7_through_theorem_4_9.tex

% BEGIN INCLUDED SGA7I_English_current_buildable_source_20260823/components/44_expose_VI_localization_and_completion_intro.tex
%%%%%%%%%%  scan idx 115  |  folio 104  %%%%%%%%%%

现在考虑概形形变理论中的局部化与完备化. 先固定记号. 设 $X_0$ 为有限型 $k$-概形, $x_0\in X_0$ 为固定点. 有自然函子
\[
\underline M_{X_0}
\longrightarrow
\underline M_{\operatorname{Spec}(\mathcal O_{X_0,x_0})}
\longrightarrow
\underline M_{\operatorname{Spec}(\widehat{\mathcal O}_{X_0,x_0})},
\]

%%%%%%%%%%  scan idx 116  |  folio 105  %%%%%%%%%%

从而有函子
\[
\widehat{\underline M}_{X_0}
\longrightarrow
\widehat{\underline M}_{\operatorname{Spec}(\mathcal O_{X_0,x_0})}
\longrightarrow
\widehat{\underline M}_{\operatorname{Spec}
 (\widehat{\mathcal O}_{X_0,x_0})}.
\]
如 \SGARefLink{sga7.sga7i.evi.n4-1.item-c}{4.1}(c) 所指出, $\widehat{\underline M}_{X_0}$ 典范等价于如下范畴: 对象是一对 $(R,\mathfrak X)$, 其中 $R\in\operatorname{Ob}(\widehat{\underline C}_\Lambda)$, $\mathfrak X$ 是平坦 $R$-进形式概形, 连同固定同构 $\mathfrak X_0\xrightarrow{\sim}X_0$.

在不致混淆时, 分别以 $(R,\mathfrak X_{(x_0)})$ 和 $(R,\mathfrak X_{(\widehat{x}_0)})$ 表示 $(R,\mathfrak X)$ 在局部化和完备化下的像. 于是, 若
\[
\mathfrak X=\varinjlim_\nu\mathfrak X_\nu,
\qquad
\mathfrak X_\nu
=\mathfrak X\times_{\operatorname{Spf}(R)}
 \operatorname{Spec}(R/\mathfrak m^{\nu+1}),
\]
则
\[
\mathfrak X_{(x_0)}
=\varinjlim_\nu
 \operatorname{Spec}(\mathcal O_{\mathfrak X_\nu,x_0})
\]
以及
\[
\mathfrak X_{(\widehat{x}_0)}
=\varinjlim_\nu
 \operatorname{Spec}(\widehat{\mathcal O}_{\mathfrak X_\nu,x_0}).
\]

% END INCLUDED SGA7I_English_current_buildable_source_20260823/components/44_expose_VI_localization_and_completion_intro.tex

% BEGIN INCLUDED SGA7I_English_current_buildable_source_20260823/components/45_expose_VI_theorem_4_10.tex
%%%%%%%%%%  scan idx 116  |  folio 105  %%%%%%%%%%

\medskip
\noindent\textbf{定理 \SGARefTarget{sga7.sga7i.evi.n4-10}{4.10}.}
设 $X_0$ 为 $k$ 上的有限型仿射概形.

\smallskip
\noindent\textbf{(a)}
若 $X_0$ 局部为完全交, 则 $\underline M_{X_0}$ 在 $\underline C_\Lambda$ 上光滑.

\smallskip
\noindent\textbf{(b)}
假设对某个固定闭点 $x\in X_0$, 开子概形 $X_0-\{x\}$ 在 $k$ 上光滑. 则两个函子
\[
\underline M_{X_0}
\longrightarrow
\underline M_{\operatorname{Spec}(\mathcal O_{X_0,x})}
\longrightarrow
\underline M_{\operatorname{Spec}(\widehat{\mathcal O}_{X_0,x})}
\]
均极小光滑. 特别地, 若 $(R,\mathfrak X)$ 为 $X_0$ 的形式遍有形变, 则 $(R,\mathfrak X_{(x)})$ 和 $(R,\mathfrak X_{(\widehat x)})$ 分别为 $\operatorname{Spec}(\mathcal O_{X_0,x})$ 和 $\operatorname{Spec}(\widehat{\mathcal O}_{X_0,x})$ 的形式遍有形变.

\medskip
\noindent\textit{证明.}
使用第 3 节交换代数上同调的记号, 并写成 $X_0=\operatorname{Spec}(A_0)$.

%%%%%%%%%%  scan idx 117  |  folio 106  %%%%%%%%%%

\smallskip
\noindent\textbf{(a)}
只需证明 $P:[\underline M_{X_0}]\longrightarrow\underline C_\Lambda$ 光滑. 设 $(R,\operatorname{Spec}(A))$ 为 $\operatorname{Spec}(A_0)$ 在 $R\in\operatorname{Ob}(\underline C_\Lambda)$ 上的形变. 给定 $\underline C_\Lambda$ 中的小扩张 $R'\longrightarrow R$, 即
\[
0\longrightarrow k\longrightarrow R'\longrightarrow R
\longrightarrow0
\]
为正合列, 我们必须将 $(R,\operatorname{Spec}(A))$ 提升到 $R'$ 上. 将系数为 $A_0$ 的 $D^i$ 的传递性序列应用于 $R'\longrightarrow R\longrightarrow A$, 得
\[
\cdots\longrightarrow D^1(A|R,A_0)
\longrightarrow D^1(A|R',A_0)
\longrightarrow D^1(R|R',k)\otimes_kA_0
\longrightarrow D^2(A|R,A_0)\longrightarrow\cdots .
\]
由于 $R'\longrightarrow R$ 为小扩张, $D^1(R|R',k)$ 在 $k$ 上为一维, 由 $0\longrightarrow k\longrightarrow R'\longrightarrow R
\longrightarrow0$ 所表示的类 $[R']$ 生成. 由于 $A$ 在 $R$ 上平坦,
\[
D^i(A|R,A_0)=D^i(A_0|k,A_0)
\qquad(i>0).
\]
特别地, 由 \SGARefLink{sga7.sga7i.evi.n3-12}{3.12}(8) 的低次描述, 因 $A_0$ 局部为完全交, 有 $D^2(A|R,A_0)=0$. 因而得到正合列
\[
\longrightarrow D^1(A_0|k,A_0)
\longrightarrow D^1(A|R',A_0)
\longrightarrow D^1(R|R',k)\otimes_kA_0
\longrightarrow0.
\]
故存在以 $A_0$ 为核的 $A$ 的 $R'$-代数扩张 $A'$, 其类映到 $[R']\otimes1$. 等价地, $(R,\operatorname{Spec}(A))$ 提升为形变 $(R',\operatorname{Spec}(A'))$.

\smallskip
\noindent\textbf{(b)}
由 \SGARefLink{sga7.sga7i.evi.n3-14}{3.14}, 有同构
\[
D^1(A_0|k,A_0)
\xrightarrow{\ \sim\ }
D^1((A_0)_x|k,(A_0)_x)
\xrightarrow{\ \sim\ }
D^1(\widehat{(A_0)_x}|k,\widehat{(A_0)_x}),
\]
从而
\[
[\underline M_{X_0}(k[\varepsilon])]
\xrightarrow{\ \sim\ }
[\underline M_{\operatorname{Spec}(\mathcal O_{X_0,x})}
 (k[\varepsilon])]
\xrightarrow{\ \sim\ }
[\underline M_{\operatorname{Spec}(\widehat{\mathcal O}_{X_0,x})}
 (k[\varepsilon])].
\]
由 \SGARefLink{sga7.sga7i.evi.n2-7}{2.7}(a), 只需证明
\[
[\underline M_{X_0}]
\longrightarrow
[\underline M_{\operatorname{Spec}(\mathcal O_{X_0,x})}]
\longrightarrow
[\underline M_{\operatorname{Spec}(\widehat{\mathcal O}_{X_0,x})}]
\]
为光滑函子.

设 $(R,\operatorname{Spec}(A))$ 为 $X_0=\operatorname{Spec}(A_0)$ 的形变, $0\longrightarrow k\longrightarrow R'\longrightarrow R
\longrightarrow0$ 为小扩张. 与上面一样, 同态 $R'\longrightarrow R\longrightarrow A$ 给出行正合的交换图

%%%%%%%%%%  scan idx 118  |  folio 107  %%%%%%%%%%

\begin{center}
\scriptsize
\begin{tikzcd}[column sep=tiny,row sep=large]
{} \arrow[r]
& D^1(A_0|k,A_0) \arrow[r] \arrow[d,"\sim"']
& D^1(A|R',A_0) \arrow[r,"p"] \arrow[d]
& D^1(R|R',k)\otimes_kA_0 \arrow[r] \arrow[d]
& D^2(A_0|k,A_0) \arrow[r] \arrow[d,"\sim"]
& \cdots \\
{} \arrow[r]
& D^1(A_0|k,(A_0)_x) \arrow[r]
& D^1(A|R',(A_0)_x) \arrow[r]
& D^1(R|R',k)\otimes_k(A_0)_x \arrow[r]
& D^2(A_0|k,(A_0)_x) \arrow[r]
& \cdots \\
{} \arrow[r]
& D^1((A_0)_x|k,(A_0)_x) \arrow[r] \arrow[u]
& D^1(A_x|R',(A_0)_x) \arrow[r,"p_x"] \arrow[u]
& D^1(R|R',k)\otimes_k(A_0)_x \arrow[r] \arrow[u]
& D^2((A_0)_x|k,(A_0)_x) \arrow[r] \arrow[u]
& \cdots
\end{tikzcd}
\end{center}
图示的同构来自如下事实: $A_0$ 除闭点 $x$ 外在 $k$ 上光滑, 从而
\[
D^i(A_0|k,(A_0)_x)
=D^i(A_0|k,A_0)_x
=D^i(A_0|k,A_0)
\qquad(i>0).
\]

假设给定 $[\underline M_{\operatorname{Spec}((A_0)_x)}]$ 中的箭头
\[
(R',\operatorname{Spec}(B))
\longrightarrow(R,\operatorname{Spec}(A_x))
\]
这相当于给定一个类 $[B]\in D^1(A_x|R',(A_0)_x)$, 满足 $p_x([B])=[R']$. 追图给出一个类 $[A']\in D^1(A|R',A_0)$, 满足 $p([A'])=[R']$ 且 $A'_x=B$. 因此 $(R',\operatorname{Spec}(A'))$ 是 $\operatorname{Spec}(A_0)$ 在 $R'$ 上的形变, 其在 $x$ 处的局部化为 $\operatorname{Spec}(B)$. 这就证明了
\[
\underline M_{X_0}\longrightarrow
\underline M_{\operatorname{Spec}(\mathcal O_{X_0,x})}
\]
光滑.

以下函子的光滑性
\[
\underline M_{\operatorname{Spec}((A_0)_x)}
\longrightarrow
\underline M_{\operatorname{Spec}(\widehat{(A_0)_x})}
\]
由类似的图以及以下等式得出:
\[
D^i((A_0)_x|k,\widehat{(A_0)_x})
=\widehat{D^i((A_0)_x|k,(A_0)_x)}
=D^i((A_0)_x|k,(A_0)_x)
\qquad(i>0),
\]
因为 $x$ 是唯一可能的非光滑点. \hfill$\square$

% END INCLUDED SGA7I_English_current_buildable_source_20260823/components/45_expose_VI_theorem_4_10.tex

% BEGIN INCLUDED SGA7I_English_current_buildable_source_20260823/components/46_expose_VI_theorem_4_11_and_remark_4_12.tex
% R3 explicit revisions VI-B-C04; see qa/FIDELITY-R3-CHANGES.json.
%%%%%%%%%%  scan idx 118  |  folio 107  %%%%%%%%%%

\medskip
\noindent\textbf{定理 \SGARefTarget{sga7.sga7i.evi.n4-11}{4.11}.}
设 $X_0$ 为 $k$ 上的有限型分离概形, 在一个有限闭子集
\[
X_0^{\mathrm{Sing}}=\{x_1,\ldots,x_r\}.
\]
以外光滑. 设 $U_0^1,\ldots,U_0^r$ 分别为 $x_1,\ldots,x_r$ 的仿射开邻域, 满足 $U_0^i\cap X_0^{\mathrm{Sing}}=\{x_i\}$. 若 $H^2(X_0,\underline D^0_{X_0})=0$, 则自然函子
\[
\underline M_{X_0}\longrightarrow
\underline M_{U_0^1}\times\cdots\times\underline M_{U_0^r}
\]
光滑.

\smallskip
\noindent\textit{证明.}
选取 $U_0^{r+1},\ldots,U_0^N$, 使对每个 $j>r$ 均有 $U_0^j\cap X_0^{\mathrm{Sing}}=\varnothing$, 将 $U_0^1,\ldots,U_0^r$ 补成 $X_0$ 的仿射开覆盖. 设 $R'\longrightarrow R$ 为 $\underline C_\Lambda$ 中的小扩张, $(R',U'^1),\ldots,(R',U'^r)$ 分别为 $U_0^1,\ldots,U_0^r$ 在 $R'$ 上的形变, 满足
\[
U'^i\otimes_{R'}R=X|_{U_0^i},\qquad 1\leq i\leq r,
\]
其中 $(R,X)$ 为 $X_0$ 在 $R$ 上的某个形变. 我们必须找到 $X_0$ 在 $R'$ 上的形变 $(R',X')$, 使得
\[
X'\otimes_{R'}R=X,
\qquad
X'|_{U_0^i}=U'^i\quad(1\leq i\leq r).
\]

%%%%%%%%%%  scan idx 119  |  folio 108  %%%%%%%%%%

对 $j>r$, 概形 $U_0^j$ 在 $k$ 上光滑; 因而由 \SGARefLink{sga7.sga7i.evi.n4-10}{4.10}(a), $X|_{U_0^j}$ 提升为 $U_0^j$ 在 $R'$ 上的形变 $(R',U'^j)$. 于是
\[
U'^j\otimes_{R'}R=X|_{U_0^j}\qquad(j>r).
\]
对每一对 $i\neq j$, 交 $U_0^i\cap U_0^j$ 仿射且在 $k$ 上光滑. 因而得到同构
\[
\tau_{ij}:U'^i|_{U_0^i\cap U_0^j}
\xrightarrow{\ \sim\ }
U'^j|_{U_0^i\cap U_0^j}
\]
使 $\tau_{ij}\otimes_{R'}R=\operatorname{id}_{X|_{U_0^i\cap U_0^j}}$. 在三重交上, 自同构
\[
\tau_{ij}\tau_{jk}\tau_{ik}^{-1}
\]
定义 $H^2(X_0,\underline D^0_{X_0})$ 中的一个上同调类. 这个类是将各 $U'^i$ 粘合为整体形变的障碍. 按假设它为零, 故各 $U'^i$ 粘合为形变 $(R',X')$, 满足 $X'\otimes_{R'}R=X$ 及 $X'|_{U_0^i}\simeq U'^i$. \hfill$\square$

\medskip
\noindent\textbf{注 \SGARefTarget{sga7.sga7i.evi.n4-12}{4.12}.}
上述证明假设 $X_0$ 分离, 但这一假设并非必要. 事实上, 固定 $X_0$ 的形变 $(R,X)$, 以及某个开覆盖各成员 $U_0^i$ 的形变 $(R',U'^i)$. $X$ 的各开子集 $V$ 在 $R'$ 上的、与已经给定的形变 $(R',U'^i)$ 相容的形变 $(R',V')$ 构成一个“gerbe” (采用 Giraud 的术语), 其结构层为 $\underline D^0_{X_0}$. 这个 gerbe 存在截面的障碍, 即存在与给定资料 $(R',U'^i)$ 相容的 $X$ 的整体形变 $(R',X')$ 的障碍, 属于 $H^2(X_0,\underline D^0_{X_0})$, 因而在定理的假设下为零.


% END INCLUDED SGA7I_English_current_buildable_source_20260823/components/46_expose_VI_theorem_4_11_and_remark_4_12.tex

% BEGIN INCLUDED SGA7I_English_current_buildable_source_20260823/components/47_expose_VI_corollary_4_13.tex
% R3 explicit revisions VI-B-C05, VI-B-C06; see qa/FIDELITY-R3-CHANGES.json.
\medskip
\noindent\textbf{推论 \SGARefTarget{sga7.sga7i.evi.n4-13}{4.13}.}
设 $X_0$ 为 $k$ 上的有限型概形, 在有限子集
\[
X_0^{\mathrm{Sing}}=\{x_1,\ldots,x_r\},
\]
以外光滑, 并假设 $H^2(X_0,\underline D^0_{X_0})=0$. 则:

\smallskip
\noindent\textbf{(a)} 函子
\[
\underline M_{X_0}\longrightarrow
\prod_{i=1}^r
\underline M_{\operatorname{Spec}(\widehat{\mathcal O}_{X_0,x_i})}
\]
光滑. 特别地,
\[
\dim\underline M_{X_0}
=\dim_k H^1(X_0,\underline D^0_{X_0})
+\sum_{i=1}^r
\dim\underline M_{\operatorname{Spec}(\widehat{\mathcal O}_{X_0,x_i})}.
\]

\smallskip
\noindent\textbf{(b)} 若此外 $X_0$ 局部为完全交, 则 $\underline M_{X_0}$ 光滑.

%%%%%%%%%%  scan idx 120  |  folio 109  %%%%%%%%%%

\smallskip
\noindent\textit{证明.}
\textbf{(a)} 第一个断言来自 \SGARefLink{sga7.sga7i.evi.n4-11}{4.11} 和 \SGARefLink{sga7.sga7i.evi.n4-10}{4.10}(b). 设 $(R,\mathfrak X)$ 以及 $(R_1,\mathfrak X_1),\ldots,(R_r,\mathfrak X_r)$ 分别为 $X_0$ 以及
\[
\operatorname{Spec}(\mathcal O_{X_0,x_1}),\ldots,
\operatorname{Spec}(\mathcal O_{X_0,x_r}),
\]
的形式遍有形变. 得到交换图
\[
\begin{tikzcd}[column sep=large,row sep=large]
\underline h_R
  \arrow[r]
  \arrow[d]
& \underline h_{R_1\widehat\otimes_\Lambda\cdots
  \widehat\otimes_\Lambda R_r}
  \arrow[d] \\
\underline M_{X_0}
  \arrow[r]
& \displaystyle\prod_{i=1}^r
  \underline M_{\operatorname{Spec}(\widehat{\mathcal O}_{X_0,x_i})}.
\end{tikzcd}
\]
复合函子
\[
\underline h_R\longrightarrow\underline M_{X_0}
\longrightarrow
\prod_{i=1}^r
\underline M_{\operatorname{Spec}(\widehat{\mathcal O}_{X_0,x_i})}
\]
光滑, 而
\[
\underline h_{R_1\widehat\otimes_\Lambda\cdots
\widehat\otimes_\Lambda R_r}
\longrightarrow
\prod_{i=1}^r
\underline M_{\operatorname{Spec}(\widehat{\mathcal O}_{X_0,x_i})}
\]
极小光滑. 由 \SGARefLink{sga7.sga7i.evi.n2-7}{2.7}(a) 得
\[
\underline h_R\longrightarrow
\underline h_{R_1\widehat\otimes_\Lambda\cdots
\widehat\otimes_\Lambda R_r}
\]
光滑; 等价地, $R$ 是 $R_1\widehat\otimes_\Lambda\cdots\widehat\otimes_\Lambda R_r$ 上的形式幂级数环.

由 \SGARefLink{sga7.sga7i.evi.n4-4}{4.4}, 序列
\[
0\longrightarrow H^1(X_0,\underline D^0_{X_0})
\longrightarrow[\underline M_{X_0}(k[\varepsilon])]
\longrightarrow
\prod_{i=1}^r
[\underline M_{\operatorname{Spec}(\widehat{\mathcal O}_{X_0,x_i})}
  (k[\varepsilon])]
\]
为正合列. 等价地,
\[
0\longrightarrow H^1(X_0,\underline D^0_{X_0})
\longrightarrow\operatorname{Hom}_{\Lambda\text{-alg}}(R,k[\varepsilon])
\longrightarrow
\operatorname{Hom}_{\Lambda\text{-alg}}
 (R_1\widehat\otimes_\Lambda\cdots\widehat\otimes_\Lambda R_r,
  k[\varepsilon])
\]
为正合列. 因此, $R$ 是 $R_1\widehat\otimes_\Lambda\cdots\widehat\otimes_\Lambda R_r$ 上的 $d$ 变元形式幂级数环, 其中
\[
d=\dim_k H^1(X_0,\underline D^0_{X_0}).
\]
因此
\begin{align*}
\dim\underline M_{X_0}
 &=\dim(k\otimes_\Lambda R) \\
 &=d+\dim\bigl(k\otimes_\Lambda
   (R_1\widehat\otimes_\Lambda\cdots\widehat\otimes_\Lambda R_r)
   \bigr) \\
 &=d+\dim\bigl((k\otimes_\Lambda R_1)
   \widehat\otimes_k\cdots\widehat\otimes_k
   (k\otimes_\Lambda R_r)\bigr) \\
 &=d+\sum_{i=1}^r\dim(k\otimes_\Lambda R_i) \\
 &=d+\sum_{i=1}^r
   \dim\underline M_{\operatorname{Spec}
   (\widehat{\mathcal O}_{X_0,x_i})}.
\end{align*}

%%%%%%%%%%  scan idx 121  |  folio 110  %%%%%%%%%%

\noindent\textbf{(b)} 此结论来自 \SGARefLink{sga7.sga7i.evi.n4-11}{4.11} 和 \SGARefLink{sga7.sga7i.evi.n4-10}{4.10}(a).
\hfill$\square$


% END INCLUDED SGA7I_English_current_buildable_source_20260823/components/47_expose_VI_corollary_4_13.tex

% BEGIN INCLUDED SGA7I_English_current_buildable_source_20260823/components/48_expose_VI_4_14_and_4_15.tex
\medskip
\noindent\textbf{\SGARefTarget{sga7.sga7i.evi.n4-14}{4.14}.}
本节最后在一个简单情形下显式构造形式遍有形变. 当 $X_0$ 为仿射完全交时, 可以立即写出这样的形变. 设
\[
A_0=k[x_1,\ldots,x_n]/(f_1,\ldots,f_r)=S/J,
\]
其中 $S=k[x_1,\ldots,x_n]$, $J=(f_1,\ldots,f_r)$, 且 $f_1,\ldots,f_r$ 为 $S$-正则序列. 有正合列
\[
\cdots\longrightarrow\operatorname{Der}_k(S,A_0)
\longrightarrow\operatorname{Hom}_{A_0}(J/J^2,A_0)
\longrightarrow D^1(A_0|k,A_0)\longrightarrow0.
\]
$A_0$-模 $J/J^2$ 是秩为 $r$ 的自由模, 由 $f_1,\ldots,f_r$ 的类生成, 因为这些元素构成一个 $S$-正则序列.

现假设 $D^1(A_0|k,A_0)$ 在 $k$ 上有限维; 例如, 可假设 $A_0$ 只有孤立非光滑点. 令
\[
d=\dim_kD^1(A_0|k,A_0),
\]
并选取映射
\[
p_i:J/J^2\longrightarrow A_0\qquad(1\leq i\leq d)
\]
使其类构成 $D^1(A_0|k,A_0)$ 的一组 $k$-基. 对 $1\leq i\leq d$ 及 $1\leq j\leq r$, 选取 $P_{ij}\in\Lambda[x_1,\ldots,x_n]$, 使
\[
P_{ij}\equiv p_i(f_j)\pmod{\mathfrak m_\Lambda+J},
\]
其中 $\mathfrak m_\Lambda$ 为 $\Lambda$ 的极大理想. 置
\[
R=\Lambda[[t_1,\ldots,t_d]],
\qquad
B=R[x_1,\ldots,x_n]/(G_1,\ldots,G_r),
\]
其中
\[
G_j=F_j-\sum_{i=1}^d t_iP_{ij},
\]
而 $F_j\in\Lambda[x_1,\ldots,x_n]$ 由将 $f_j$ 的系数从 $k$ 提升到 $\Lambda$ 而得.

代数 $R[x_1,\ldots,x_n]$ 在 $R$ 上平坦, 而 $G_1,\ldots,G_r$ 模 $\mathfrak m_R$ 的像构成 $k\otimes_RR[x_1,\ldots,x_n]$ 中的正则序列 $f_1,\ldots,f_r$. 因而 $B$ 在 $R$ 上平坦, 且 $k\otimes_RB=A_0$. 按构造, 映射
\[
\operatorname{Hom}_{\Lambda\text{-alg}}(R,k[\varepsilon])
\longrightarrow D^1(A_0|k,A_0)
=[\underline M_{\operatorname{Spec}(A_0)}(k[\varepsilon])]
\]
为双射. 因此, 若

%%%%%%%%%%  scan idx 122  |  folio 111  %%%%%%%%%%

\[
\operatorname{Spf}(B)
=\varinjlim_\nu
\operatorname{Spec}\bigl((R/\mathfrak m_R^{\nu+1})\otimes_RB\bigr),
\]
则 $(R,\operatorname{Spf}(B))$ 为 $\widehat{\underline M}_{\operatorname{Spec}(A_0)}$ 的一个对象. 其诱导的函子
\[
\underline h_R\longrightarrow\underline M_{\operatorname{Spec}(A_0)}
\]
由 \SGARefLink{sga7.sga7i.evi.n1-6}{1.6} 是光滑的, 因为 $R$ 为 $\Lambda$ 上的形式幂级数环. 换言之, $(R,\operatorname{Spf}(B))$ 为 $\underline M_{\operatorname{Spec}(A_0)}$ 的形式遍有对象.

由 \SGARefLink{sga7.sga7i.evi.n4-10}{4.10}(b), 若改为置
\[
B=R[[x_1,\ldots,x_n]]/(G_1,\ldots,G_r),
\]
则 $(R,\operatorname{Spf}(B))$ 是 $\operatorname{Spec}(A_0)$ 的形式遍有形变, 其中
\[
A_0=k[[x_1,\ldots,x_n]]/(f_1,\ldots,f_r).
\]

\medskip
\noindent\textbf{注 \SGARefTarget{sga7.sga7i.evi.n4-15}{4.15}.}
设 $X_0$ 为 $k$ 上的有限型概形, 仅有孤立非光滑点, $(R,\mathfrak X)$ 为 $X_0$ 的形式遍有形变. 上述构造表明, 当 $X_0$ 仿射且局部为完全交时, $\mathfrak X$ 是某个有限型 $R$-概形的形式化. 只要 $H^2(X_0,\underline D^0_{X_0})=0$, 这一事实也容易推广到非仿射情形.


% END INCLUDED SGA7I_English_current_buildable_source_20260823/components/48_expose_VI_4_14_and_4_15.tex

% BEGIN INCLUDED SGA7I_English_current_buildable_source_20260823/components/49_expose_VI_remark_4_16.tex
\medskip
\noindent\textbf{注 \SGARefTarget{sga7.sga7i.evi.n4-16}{4.16}.}
相对余切复形的整体理论推广并简化了诸如 \SGARefLink{sga7.sga7i.evi.n4-10}{4.10} 的结果; 这里将复形截断为长度一的复形, 但这已足以处理形变理论中的许多问题. 设 $S$ 为概形, $S_0\subset S'\subset S$ 为分别由理想 $J\supset I$ 定义的子概形, 且 $JI=0$. 设 $X'$ 在 $S'$ 上平坦, 并以
\[
L_\bullet=L^{X_0/S_0}
\]
表示 $X_0$ 在 $S_0$ 上的相对余切复形, 将其视为 $D(X_0,\mathcal O_{X_0})$ 的对象. 设 $\underline F$ 为延拓 $X'/S'$ 的平坦 $S$-概形 $X$ 所成的范畴; 因而一个对象带有 $S'$-同构
\[
X\times_SS'\xrightarrow{\ \sim\ }X'.
\]
则:

\smallskip
\noindent\textbf{(a)} 对 $\underline F$ 的一个对象 $X$, 其在 $X'$ 上诱导恒等的自同构群典范同构于
\[
\operatorname{Ext}^0(X_0;L_\bullet,I_{X_0}),
\]
其中 $I_{X_0}$ 是 $\mathcal O_{S_0}$-模 $I$ 在 $X_0$ 上的逆像.

%%%%%%%%%%  scan idx 123  |  folio 112  %%%%%%%%%%

\smallskip
\noindent\textbf{(b)} $\underline F$ 的对象同构类集合或者为空, 或者是 $\operatorname{Ext}^1(X_0;L_\bullet,I_{X_0})$ 下的挠子 (torsor, 即主齐性空间).

\smallskip
\noindent\textbf{(c)} 存在典范元素
\[
c\in\operatorname{Ext}^2(X_0;L_\bullet,I_{X_0})
\]
其为零是 $\underline F$ 非空的充要条件.

为研究这些群, 往往方便写成
\[
\operatorname{Ext}^i(X_0;L_\bullet,I_{X_0})
=H^i\bigl(X_0,R\underline{\operatorname{Hom}}
  (L_\bullet,I_{X_0})\bigr).
\]
$R\underline{\operatorname{Hom}}(L_\bullet,I_{X_0})$ 的上同调层恰好是第 3 节引入的层 $\underline D^i_{X_0}(I_{X_0})$. 对 $i=0,1,2$, 它们分别可视为 $\underline F$ 的对象 $X$ 的自同构层 (这些自同构在 $X'$ 上诱导恒等)、度量平坦形变问题局部解的同构类不确定性的层, 以及局部解存在性的障碍层.

标准谱序列为
\[
E_2^{p,q}
=H^p\bigl(X_0,\underline D^q_{X_0}(I_{X_0})\bigr)
\Longrightarrow
\operatorname{Ext}^{p+q}(X_0;L_\bullet,I_{X_0}).
\]
它给出熟知的同构
\begin{align*}
\operatorname{Aut}(X/X')
&\simeq H^0\bigl(X_0,\underline D^0_{X_0}(I_{X_0})\bigr)\\
&\simeq H^0\bigl(X_0,
  \underline{\operatorname{Hom}}
  (\Omega^1_{X_0/S_0},I_{X_0})\bigr)\\
&=\operatorname{Hom}(\Omega^1_{X_0/S_0},I_{X_0}),
\end{align*}
以及低次正合列
\[
0\longrightarrow H^1(X_0,\underline D^0_{X_0}(I_{X_0}))
\longrightarrow\operatorname{Ext}^1
\longrightarrow H^0(X_0,\underline D^1_{X_0}(I_{X_0}))
\longrightarrow H^2(X_0,\underline D^0_{X_0}(I_{X_0}))
\longrightarrow\operatorname{Ext}^2.
\]
特别地, 若 $H^2(X_0,\underline D^0_{X_0}(I_{X_0}))=0$, 且 $X'/S'$ 的局部平坦形变存在, 则相容的局部形变可粘合为整体平坦形变. 更一般地, 对局部平坦形变同构类所成的层的一个指定截面, \SGARefLink{sga7.sga7i.evi.n4-10}{4.10} 的论证将相容整体形变的障碍置于 $H^2(X_0,\underline D^0_{X_0}(I_{X_0}))$ 中.

%%%%%%%%%%  scan idx 124  |  folio 113  %%%%%%%%%%

同一谱序列表明, 只要
\begin{equation}\tag{*}
H^0(X_0,\underline D^2_{X_0})=0,
\qquad
H^1(X_0,\underline D^1_{X_0})=0,
\qquad
H^2(X_0,\underline D^0_{X_0})=0.
\end{equation}
就有 $\operatorname{Ext}^2$ 为零. 当 $X_0$ 为 $S_0$ 上的相对完全交时, 第一个等式总成立, 因为此时对每个 $\mathcal O_{X_0}$-模 $M$ 都有 $\underline D^2_{X_0}(M)=0$. 当 $\underline D^1_{X_0}$ 的支集离散时, 第二个等式成立; 当 $X_0$ 除一个离散子集外在 $S_0$ 上光滑时便是这种情形. 若 $(*)$ 中的三个等式全成立, 则由 (c), $\underline F$ 非空, 这推广了 \SGARefLink{sga7.sga7i.evi.n4-13}{4.13}(b).

还有一种几何解释, 它只需要整体化的函子 $\underline D^i_{X_0}$, 而不需要相对余切复形的整体理论. 将 $X'/S'$ 平坦地延拓到 $S$ 上的三个相继的障碍, 分别位于 $(*)$ 所列的三个群中. 局部理论定义第一个障碍
\[
c_1\in H^0(X_0,\underline D^2_{X_0})
\]
即局部解存在性的障碍. 若 $c_1=0$, 则局部解同构类所成的层 $\underline F$ 典范地为 $\underline D^1_{X_0}$ 下的挠子. 它存在截面的障碍为
\[
c_2\in H^1(X_0,\underline D^1_{X_0});
\]
这样的截面是局部解的同构类的一个相容系统. 若 $c_2=0$, 选取 $\underline F$ 的一个截面. 此时, 存在实现所选局部同构类的整体解的障碍为
\[
c_3\in H^2(X_0,\underline D^0_{X_0}).
\]
用 $\underline D^1_{X_0}$ 的一个截面改变 $\underline F$ 的截面, 给出映射
\[
d:H^0(X_0,\underline D^1_{X_0})
\longrightarrow H^2(X_0,\underline D^0_{X_0}).
\]
解存在, 当且仅当 $c_1=0$、$c_2=0$, 且 $c_3$ 在 $\operatorname{Coker}d$ 中的类为零 (每个 $c_i$ 都只在此前各障碍为零之后才有定义).

%%%%%%%%%%  scan idx 125  |  folio 114  %%%%%%%%%%

这三个障碍群
\[
H^0(X_0,\underline D^2_{X_0}),
\qquad
H^1(X_0,\underline D^1_{X_0}),
\qquad
H^2(X_0,\underline D^0_{X_0})/\operatorname{Im}d
\]
是上述谱序列的起始项 $E_2^{0,2}$、$E_2^{1,1}$ 和 $E_3^{2,0}$; 在相差一个符号的意义下, $d$ 就是它的微分
\[
d_2:E_2^{0,1}\longrightarrow E_2^{2,0}.
\]


% END INCLUDED SGA7I_English_current_buildable_source_20260823/components/49_expose_VI_remark_4_16.tex

% BEGIN INCLUDED SGA7I_English_current_buildable_source_20260823/components/50_expose_VI_section_5_opening_through_5_2.tex
\bigskip
\noindent\textbf{5. 形式 Jacobian 子概形}

\medskip
设 $X_0$ 为域 $k$ 上的有限型概形, $(R,\mathfrak X)$ 为 $X_0$ 的形式遍有形变. 空间 $\mathfrak X$ 是一个 $R$-进形式概形, 并非通常的 $R$-概形. 本节旨在阐明“$\mathfrak X$ 在 $R$ 上的非光滑轨迹”的含义.

首先回顾有限型模的 Fitting 理想. 设 $A$ 为交换环, $E$ 为有限型 $A$-模. 选取一个表示
\[
K\longrightarrow F\longrightarrow E\longrightarrow0,
\]
其中 $F$ 局部自由, 恒秩为 $n$. 对 $p=0,1,2,\ldots$, 定义
\[
I_A^p(E)=\operatorname{Im}\left(
 \bigwedge^{n-p}K\otimes_A\bigwedge^{n-p}F^*
 \longrightarrow A\right).
\]
这些理想与所选的表示无关. 它们是 $A$-模 $E$ 的不变量, 称为其 \emph{Fitting 理想}. 下列性质可直接由定义得出.

\medskip
\noindent\textbf{\SGARefTarget{sga7.sga7i.evi.n5-1}{5.1}.}
\textbf{(a)} 对每个环同态 $A\longrightarrow B$, 典范映射
\[
B\otimes_A A/I_A^p(E)
\xrightarrow{\ \sim\ }
B/I_B^p(B\otimes_AE)
\]
对每个 $p$ 都是同构. 特别地,
\[
I_{S^{-1}A}^p(S^{-1}E)=S^{-1}I_A^p(E)
\]
对每个乘法子集 $S\subset A$ 都成立.

%%%%%%%%%%  scan idx 126  |  folio 115  %%%%%%%%%%

\smallskip
\noindent\textbf{(b)} 有包含关系
\[
I_A^0(E)\subset I_A^1(E)\subset I_A^2(E)\subset\cdots.
\]
对每个 $x\in\operatorname{Spec}(A)$,
\[
I_A^p(E)_x=A_x
\qquad\text{对所有 }p\geq\dim_{\kappa(x)}E(x),
\]
其中 $E(x)=\kappa(x)\otimes_AE$.

\smallskip
\noindent\textbf{(c)} 点 $x\in\operatorname{Spec}(A)-\operatorname{Supp}(A/I_A^p(E))$, 当且仅当对任意一个 (等价地, 某一个) 表示
\[
K\longrightarrow F\longrightarrow E\longrightarrow0,
\qquad F\simeq A^n,
\]
模 $K_x$ 包含一个同构于 $A_x^{n-p}$ 的子模, 它是 $F_x$ 的直和项. 特别地,
\[
\operatorname{Supp}(A/I_A^0(E))=\operatorname{Supp}(E).
\]

\smallskip
\noindent\textbf{(d)} 对任意两个有限型 $A$-模,
\[
I_A^r(E_1\oplus E_2)
=\sum_{p+q=r}I_A^p(E_1)I_A^q(E_2).
\]

\smallskip
\noindent\textbf{(e)} 若 $A$ 为 Noether 进环, 则
\[
A/I_A^p(E)=\varprojlim_\nu A_\nu/I_{A_\nu}^p(E_\nu),
\]
其中 $A=\varprojlim_\nu A_\nu$, $E_\nu=E\otimes_AA_\nu$.

性质 \SGARefLink{sga7.sga7i.evi.n5-1}{5.1}(a) 使这一定义整体化. 若 $X$ 为概形, $\underline E$ 为有限表示的拟凝聚 $\mathcal O_X$-模, 则存在拟凝聚理想层 $\underline I_X^p(\underline E)$, $p\geq0$, 它们在每个仿射开子集 $U\subset X$ 上由下式刻画:
\[
\Gamma(U,\underline I_X^p(\underline E))
\xrightarrow{\ \sim\ }
I_{\Gamma(U,\mathcal O_X)}^p(\Gamma(U,\underline E)).
\]

对概形态射 $X\longrightarrow Y$, 相对 K\"ahler 微分模 $\Omega_{X|Y}=\Omega^1_{X|Y}$ 是拟凝聚的. 若此态射本质有限表示, 则 $\Omega_{X|Y}$ 为有限表示, 因而理想层 $\underline I_X^p(\Omega_{X|Y})$ 有定义.

\medskip
\noindent\textbf{定义 \SGARefTarget{sga7.sga7i.evi.n5-2}{5.2}.}
设 $X\longrightarrow Y$ 为本质有限表示态射. 对每个整数 $p\geq0$, 定义 $J_X^p(X|Y)$ 为 $X$ 的闭子概形

%%%%%%%%%%  scan idx 127  |  folio 116  %%%%%%%%%%

\noindent 它由理想层
\[
\underline I_X^p(X|Y)=\underline I_X^p(\Omega_{X|Y}).
\]
定义, 称为 $X$ 在 $Y$ 上的第 $p$ 个 \emph{Jacobian 子概形}. 对每个 $x\in X$, 置
\[
X_x=\operatorname{Spec}(\mathcal O_{X,x}).
\]


% END INCLUDED SGA7I_English_current_buildable_source_20260823/components/50_expose_VI_section_5_opening_through_5_2.tex

% BEGIN INCLUDED SGA7I_English_current_buildable_source_20260823/components/51_expose_VI_theorem_5_3.tex
\medskip
\noindent\textbf{定理 \SGARefTarget{sga7.sga7i.evi.n5-3}{5.3}.}
设 $X\longrightarrow Y$ 为本质有限表示态射. 则:

\smallskip
\noindent\textbf{(1)} 对每个基变换 $Y'\longrightarrow Y$, 典范映射
\[
J_{X'}^p(X'|Y')\xrightarrow{\ \sim\ }
J_X^p(X|Y)\times_YY'
\]
对每个 $p$ 都是同构, 其中 $X'=X\times_YY'$.

\smallskip
\noindent\textbf{(2)} 对每个 $x\in X$, 存在典范同构
\[
J_{X_x}^p(X_x|Y)\xrightarrow{\ \sim\ }J_X^p(X|Y)_x.
\]

\smallskip
\noindent\textbf{(3)} 有包含关系
\[
J_X^0(X|Y)\supset J_X^1(X|Y)\supset J_X^2(X|Y)\supset\cdots,
\]
以及
\[
x\notin J_X^p(X|Y)
\qquad\text{对每个 }p\geq
\dim_{\kappa(x)}\Omega_{X|Y}(x).
\]

\smallskip
\noindent\textbf{(4)} 点 $x$ 不属于 $J_X^p(X|Y)$, 当且仅当它有一个开邻域 $U$, 可通过闭浸入 $i:U\hookrightarrow U'$ 嵌入某个 $Y$-概形 $U'$, 使得 $U'$ 在 $x'=i(x)$ 处在 $Y$ 上光滑, 且
\[
\dim_{x'}U'_{f(x)}=p.
\]

\smallskip
\noindent\textbf{(5)} 对投影 $\pi_i:X_1\times_YX_2\longrightarrow X_i$,
\[
\underline I_{X_1\times_YX_2}^r(X_1\times_YX_2|Y)
=\sum_{p+q=r}
 \pi_1^*\underline I_{X_1}^p(X_1|Y)\,
 \pi_2^*\underline I_{X_2}^q(X_2|Y).
\]

\smallskip
\noindent\textit{证明.}
断言 (3) 来自 \SGARefLink{sga7.sga7i.evi.n5-1}{5.1}(b). 断言 (1) 来自 \SGARefLink{sga7.sga7i.evi.n5-1}{5.1}(a) 及典范同构
\[
P^*\Omega_{X|Y}\xrightarrow{\ \sim\ }\Omega_{X'|Y'},
\qquad P:X\times_YY'\longrightarrow X.
\]
同样, (2) 来自 \SGARefLink{sga7.sga7i.evi.n5-1}{5.1}(a) 以及
\[
\mathcal O_{X,x}\otimes_{\mathcal O_X}\Omega_{X|Y}
\xrightarrow{\ \sim\ }\Omega_{X_x|Y}.
\]
最后, (5) 来自 \SGARefLink{sga7.sga7i.evi.n5-1}{5.1}(d) 以及
\[
\pi_1^*\Omega_{X_1|Y}\oplus\pi_2^*\Omega_{X_2|Y}
\xrightarrow{\ \sim\ }
\Omega_{X_1\times_YX_2|Y}.
\]

对于 (4), 首先假设有图
\[
\begin{tikzcd}[column sep=huge,row sep=large]
U \arrow[r,hook,"i"] \arrow[d,"f"']
  & U' \arrow[dl] \\
Y & {}
\end{tikzcd}
\]
其中 $U'$ 在 $x'=i(x)$ 处在 $Y$ 上光滑, 且 $\dim_{x'}U'_{f(x)}=p$. 则
\[
\dim_{\kappa(x')}\Omega_{U'|Y}(x')=p,
\]
且序列
\[
i^*\Omega_{U'|Y}\longrightarrow\Omega_{U|Y}\longrightarrow0
\]
为正合列. 从而 $\dim_{\kappa(x)}\Omega_{U|Y}(x)\leq p$, 故由 (3), $x\notin J_U^p(U|Y)$. 由于
\[
J_U^p(U|Y)=J_X^p(X|Y)|_U,
\]
得到 $x\notin J_X^p(X|Y)$.

%%%%%%%%%%  scan idx 128  |  folio 117  %%%%%%%%%%

反之, 假设 $x\notin J_X^p(X|Y)$. 问题是局部的, 因而可写成 $X=\operatorname{Spec}(B)$ 和 $Y=\operatorname{Spec}(A)$. 选取一个 $A$-表示
\[
0\longrightarrow J\longrightarrow P\longrightarrow B
\longrightarrow0,
\qquad P=A[x_1,\ldots,x_n].
\]
有正合列
\[
J/J^2\xrightarrow{d}B\otimes_P\Omega_{P|A}
\longrightarrow\Omega_{B|A}\longrightarrow0,
\]
且 $B\otimes_P\Omega_{P|A}$ 是 $B$ 上的秩 $n$ 自由模. 因此
\[
I_B^p(\Omega_{B|A})
=\operatorname{Im}\left(
 \bigwedge^{n-p}(J/J^2)\otimes_B
 (B\otimes_P\bigwedge^{n-p}\Omega_{P|A})^*
 \longrightarrow B\right).
\]
因此, $x\notin J_X^p(X|Y)$ 当且仅当存在 $f_1,\ldots,f_{n-p}\in J$, 并可从变元 $x_1,\ldots,x_n$ 中选出 $n-p$ 个, 使相应的 Jacobian 子式
\[
\det\left(\frac{\partial f_i}{\partial x_j}\right)(x)
\]
非零. 等价地,
\[
\operatorname{Spec}\bigl(P/(f_1,\ldots,f_{n-p})\bigr)
\longrightarrow\operatorname{Spec}(A)
\]
在 $x$ 的像 $x'$ 处光滑. 取
\[
U'=\operatorname{Spec}\bigl(P/(f_1,\ldots,f_{n-p})\bigr)
\]
即局部给出所需的闭浸入. \hfill$\square$


% END INCLUDED SGA7I_English_current_buildable_source_20260823/components/51_expose_VI_theorem_5_3.tex

% BEGIN INCLUDED SGA7I_English_current_buildable_source_20260823/components/52_expose_VI_theorem_5_4.tex
\medskip
\noindent\textbf{定理 \SGARefTarget{sga7.sga7i.evi.n5-4}{5.4}.}
设 $X\xrightarrow{f}Y$ 为本质有限表示态射. 假设存在固定整数 $n\geq0$, 使下列条件之一成立:

\smallskip
\noindent\textbf{(1)} $f$ 平坦, 且对每个 $x\in X$ 都有 $\dim_xX_{f(x)}=n$;

\smallskip
\noindent\textbf{(2)} $f$ 为虚维数 $n$ 的相对完全交 (SGA 6, VIII, 1.9).

\smallskip
则 $x\notin J_X^n(X|Y)$, 当且仅当 $f$ 在 $x$ 处光滑.

\smallskip
\noindent\textit{证明.}
由 \SGARefLink{sga7.sga7i.evi.n5-3}{5.3}(4), $x\notin J_X^n(X|Y)$ 当且仅当有图
\[
\begin{tikzcd}[column sep=huge,row sep=large]
U \arrow[r,hook,"i"] \arrow[d,"f"']
  & U' \arrow[dl,"f'"] \\
Y & {}
\end{tikzcd}
\]
使得 $f'$ 在 $x'$ 处光滑, $\dim_{x'}U'_{f(x)}=n$, 而 $i$ 为闭浸入, $x'=i(x)$. 因而只需证明, 在 (1) 或 (2) 中任一条件下, $i$ 都在 $x$ 处为同构.

置
\[
A=\mathcal O_{Y,f(x)},\qquad
B'=\mathcal O_{U',x'},\qquad B=\mathcal O_{U,x},
\]
并设
\[
0\longrightarrow I\longrightarrow B'\longrightarrow B
\longrightarrow0
\]
为正合列. 必须证明 $I=0$.

假设 (1) 成立. 由于 $B$ 在 $A$ 上平坦, 序列
\[
0\longrightarrow I/\mathfrak mI
\longrightarrow B'/\mathfrak mB'
\longrightarrow B/\mathfrak mB\longrightarrow0
\]
为正合列, 其中 $\mathfrak m$ 为 $A$ 的极大理想. 按假设,
\[
\dim_xU_{f(x)}=\dim_{x'}U'_{f(x)},
\]
从而 $\dim(B'/\mathfrak mB')=\dim(B/\mathfrak mB)$. 由于 $B'/\mathfrak mB'$ 为正则局部环, 得到 $I/\mathfrak mI=0$, 因而 $I=0$.

%%%%%%%%%%  scan idx 129  |  folio 118  %%%%%%%%%%

现在假设 (2) 成立, 并考虑正合列
\[
I/I^2\xrightarrow{d}B\otimes_{B'}\Omega_{B'|A}
\longrightarrow\Omega_{B|A}\longrightarrow0.
\]
由于 $X$ 为 $Y$ 上的相对完全交, 且 $U'\longrightarrow Y$ 在 $x'$ 处光滑, 映射
\[
d:I/I^2\longrightarrow B\otimes_{B'}\Omega_{B'|A}
\]
在 $x$ 处为单射, 且 $I/I^2$ 与 $B\otimes_{B'}\Omega_{B'|A}$ 均在 $x$ 处自由. 因而由相对完全交虚维数的定义, 必有
\[
n=\dim_{\kappa(x)}\Omega_{B'|A}(x')
  -\dim_{\kappa(x)}(I/I^2)(x)
 =n-\dim_{\kappa(x)}(I/I^2)(x).
\]
故 $I/I^2=0$, 从而 $I=0$. \hfill$\square$

% END INCLUDED SGA7I_English_current_buildable_source_20260823/components/52_expose_VI_theorem_5_4.tex

% BEGIN INCLUDED SGA7I_English_current_buildable_source_20260823/components/53_expose_VI_5_5_through_5_6.tex
% R3 explicit revisions VI-B-C08; see qa/FIDELITY-R3-CHANGES.json.
\medskip
\noindent\textbf{定义 \SGARefTarget{sga7.sga7i.evi.n5-5}{5.5}.}
设 $X\xrightarrow{f}Y$ 为本质有限表示态射. 若存在固定整数 $n$ 满足 \SGARefLink{sga7.sga7i.evi.n5-4}{5.4} 的 (1) 或 (2), 置
\[
\underline I_X(X|Y)=\underline I_X^n(X|Y),
\qquad
J_X(X|Y)=J_X^n(X|Y).
\]

%%%%%%%%%%  scan idx 130  |  folio 119  %%%%%%%%%%

这一概念自然推广到形式概形. 设 $\mathfrak X$ 为 诺特形式概形, $\mathcal E$ 为凝聚 $\mathfrak X$-模. 对每个仿射开子集 $U\subset\mathfrak X$, 定义
\[
\underline I_{\mathfrak X}^p(\mathcal E)(U)
=I_{\mathfrak X(U)}^p\bigl(\mathcal E(U)\bigr).
\]
若 $\mathfrak X=\varinjlim_\nu X_\nu$, 则 $\mathcal E=\varprojlim_\nu\mathcal E_\nu$, 其中
\[
\mathcal E_\nu
=\mathcal E\otimes_{\mathcal O_{\mathfrak X}}\mathcal O_{X_\nu},
\]
且由 Noether 性得
\[
\underline I_{\mathfrak X}^p(\mathcal E)
=\varprojlim_\nu\underline I_{X_\nu}^p(\mathcal E_\nu).
\]

现设 $S$ 为 诺特形式概形, $\mathfrak X$ 为本质有限型的 $S$-进形式概形. 设 $\underline I_{\mathfrak X}$ 为定义对角闭浸入
\[
\Delta:\mathfrak X\longrightarrow
\mathfrak X\times_S\mathfrak X
\]
的理想层 (纤维积在形式概形范畴中取), 因而
\[
0\longrightarrow\underline I_{\mathfrak X}
\longrightarrow\mathcal O_{\mathfrak X\times_S\mathfrak X}
\longrightarrow\mathcal O_{\mathfrak X}\longrightarrow0
\]
为正合列. 定义
\[
\Omega_{\mathfrak X|S}
=\underline I_{\mathfrak X}/\underline I_{\mathfrak X}^{2}.
\]
写成
\[
\mathfrak X=\varinjlim_\nu X_\nu,
\qquad X_\nu=\mathfrak X\times_SS_\nu,
\]
有
\[
\Omega_{\mathfrak X|S}
\simeq\varprojlim_\nu\Omega_{X_\nu|S_\nu},
\]
故 $\Omega_{\mathfrak X|S}$ 为凝聚 $\mathfrak X$-模. 对每个有限型形式 $S$-概形 $S'$, 存在典范同构
\[
\Omega_{\mathfrak X\times_SS'|S'}
\simeq
\mathcal O_{S'}\otimes_{\mathcal O_S}\Omega_{\mathfrak X|S}.
\]

\medskip
\noindent\textbf{定义 \SGARefTarget{sga7.sga7i.evi.n5-6}{5.6}.}
设 $S$ 为 诺特形式概形, $\mathfrak X$ 为本质有限型的 $S$-进形式概形. 定义
\[
\underline I_{\mathfrak X}^p(\mathfrak X|S)
=\underline I_{\mathfrak X}^p(\Omega_{\mathfrak X|S}),
\]
并设 $J_{\mathfrak X}^p(\mathfrak X|S)$ 为此理想层定义的 $\mathfrak X$ 的闭形式子概形.

由于 $\Omega_{\mathfrak X|S}=\varprojlim_\nu\Omega_{X_\nu|S_\nu}$, 有
\[
\underline I_{\mathfrak X}^p(\mathfrak X|S)
=\varprojlim_\nu\underline I_{X_\nu}^p(X_\nu|S_\nu)
\]
以及
\[
J_{\mathfrak X}^p(\mathfrak X|S)
=\varinjlim_\nu J_{X_\nu}^p(X_\nu|S_\nu).
\]

%%%%%%%%%%  scan idx 131  |  folio 120  %%%%%%%%%%

形式概形 $\mathfrak X_{(x)}$ 仍为本质有限型的 $S$-进形式概形, 且有交换图
\[
\begin{tikzcd}[column sep=huge,row sep=large]
\mathfrak X_{(\widehat{x})}
  \arrow[rr,"i_{(\widehat{x})}"] \arrow[dr]
  && \mathfrak X \\
& \mathfrak X_{(x)} \arrow[ur,"i_{(x)}"'] & {}
\end{tikzcd}
\]
% Direct 1800-dpi source check: the upper-left subscript and its arrow
% label are both \widehat{x}; the lower subscript and label are plain x.
若 $x$ 为 $\mathfrak X$ 的孤立点, 则 $i_{(x)}$ 与 $i_{(\widehat{x})}$ 都同构到 $\mathfrak X$ 中包含 $x$ 的连通分支上.

特别地, 若
\[
|\mathfrak X|=\{x_1,\ldots,x_m\}
\]
为有限离散空间, 则有典范等同
\[
\mathfrak X=\coprod_{1\leq i\leq m}\mathfrak X_{(x_i)}
\qquad\text{(不交并)}.
\]

% END INCLUDED SGA7I_English_current_buildable_source_20260823/components/53_expose_VI_5_5_through_5_6.tex

% BEGIN INCLUDED SGA7I_English_current_buildable_source_20260823/components/54_expose_VI_corollary_5_7.tex
\medskip
\noindent\textbf{推论 \SGARefTarget{sga7.sga7i.evi.n5-7}{5.7}.}
设 $S$ 为 诺特形式概形, $\mathfrak X$ 为本质有限型的 $S$-进形式概形. 则:

\smallskip
\noindent\textbf{(1)} $J_{\mathfrak X}^p(\mathfrak X|S)$ 的形成与基变换可交换. 更确切地, 若
\[
\mathfrak X'=\mathfrak X\times_SS',
\]
则
\[
J_{\mathfrak X'}^p(\mathfrak X'|S')
=J_{\mathfrak X}^p(\mathfrak X|S)\times_SS'.
\]

\smallskip
\noindent\textbf{(2)} $J_{\mathfrak X}^p(\mathfrak X|S)$ 的连通分支与 $J_{X_0}^p(X_0|S_0)$ 的连通分支一一对应. 对每个
\[
x\in\bigl|J_{X_0}^p(X_0|S_0)\bigr|
\]
存在典范同构
\[
J_{\mathfrak X}^p(\mathfrak X|S)_{(x)}
\simeq
J_{\mathfrak X_{(x)}}^p(\mathfrak X_{(x)}|S).
\]
特别地, 若
\[
\bigl|J_{X_0}^p(X_0|S_0)\bigr|=\{x_1,\ldots,x_m\}
\]
为离散集, 则
\[
J_{\mathfrak X}^p(\mathfrak X|S)
=\coprod_{1\leq i\leq m}
 J_{\mathfrak X_{(x_i)}}^p(\mathfrak X_{(x_i)}|S)
\qquad\text{(不交并)}.
\]

%%%%%%%%%%  scan idx 132  |  folio 121  %%%%%%%%%%

\smallskip
\noindent\textbf{(3)} 设 $S$ 为通常的 诺特概形, $X$ 为本质有限型 $S$-概形. 给定闭子概形 $S_0\subset S$, 设 $\widehat S$ 为 $S$ 沿 $S_0$ 的形式完备化. 若
\[
\widehat X=X\times_S\widehat S
\]
为 $X$ 沿 $X_0=X\times_SS_0$ 的形式完备化, 则
\[
J_{\widehat X}^p(\widehat X|\widehat S)
=J_X^p(X|S)\times_S\widehat S.
\]
等价地, 它是 $J_X^p(X|S)$ 沿 $J_{X_0}^p(X_0|S_0)$ 的形式完备化.

\smallskip
\noindent\textbf{(4)} 假设固定整数 $n\geq0$ 对每个 $X_\nu\longrightarrow S_\nu$ 均满足条件 \SGARefLink{sga7.sga7i.evi.n5-4}{5.4}(1) 或 \SGARefLink{sga7.sga7i.evi.n5-4}{5.4}(2). 则
\[
x\notin J_{\mathfrak X}^n(\mathfrak X|S)
\]
当且仅当 $\mathfrak X$ 在 $x$ 处在 $S$ 上形式光滑.

\smallskip
\noindent\textit{证明.}
断言 (1) 来自
\[
\Omega_{\mathfrak X'|S'}
\simeq
\mathcal O_{S'}\widehat\otimes_{\mathcal O_S}
\Omega_{\mathfrak X|S},
\]
而 (2) 来自
\[
\mathcal O_{\mathfrak X_{(x)}}
\widehat\otimes_{\mathcal O_{\mathfrak X}}
\Omega_{\mathfrak X|S}
\simeq\Omega_{\mathfrak X_{(x)}|S}.
\]
对于 (3),
\[
\begin{aligned}
J_{\widehat X}^p(\widehat X|\widehat S)
&=\varinjlim_\nu
 J_{X\times_SS_\nu}^p(X\times_SS_\nu|S_\nu) \\
&=\varinjlim_\nu\bigl(J_X^p(X|S)\times_SS_\nu\bigr)
 =J_X^p(X|S)\times_S\widehat S.
\end{aligned}
\]
最后, 由于 $\mathfrak X$ 为 $S$-进形式概形,
\[
\begin{aligned}
\mathfrak X\to S\text{ 在以下点处形式光滑: }x
&\Longleftrightarrow
 X_\nu\to S_\nu\text{ 在以下点处形式光滑: }x
 \text{ 对每个 }\nu \\
&\Longleftrightarrow
 x\notin J_{X_\nu}^n(X_\nu|S_\nu)
 \text{ 对每个 }\nu \\
&\Longleftrightarrow
 x\notin\varinjlim_\nu J_{X_\nu}^n(X_\nu|S_\nu)
 =J_{\mathfrak X}^n(\mathfrak X|S).
\end{aligned}
\]
\hfill$\square$

% END INCLUDED SGA7I_English_current_buildable_source_20260823/components/54_expose_VI_corollary_5_7.tex

% BEGIN INCLUDED SGA7I_English_current_buildable_source_20260823/components/55_expose_VI_remark_5_8.tex
\medskip
\noindent\textbf{注 \SGARefTarget{sga7.sga7i.evi.n5-8}{5.8}.}
设 $\mathfrak X$ 为本质有限型的 $S$-进形式概形. 对点 $x\in\mathfrak X$, 另考虑 $S$-进形式概形
\[
\mathfrak X_{(\widehat{x})}
=\varinjlim_\nu
 \operatorname{Spec}\bigl(\widehat{\mathcal O_{X_\nu,x}}\bigr).
\]
有自然态射
\[
\mathfrak X_{(\widehat{x})}
\longrightarrow\mathfrak X_{(x)}
\longrightarrow\mathfrak X.
\]
一般地, $\widehat{\mathcal O_{X_0,x}}$ 在 $S_{(y)}$ 上并非本质有限型, 其中 $y$ 为 $x$ 在 $\mathfrak X\longrightarrow S$ 下的像.

为简便起见, 令
\[
A_\nu=\mathcal O_{X_\nu,x},
\qquad R_\nu=\mathcal O_{S_\nu,y}.
\]
自然映射
\[
\Omega_{A_\nu|R_\nu}
\longrightarrow\widehat\Omega_{A_\nu|R_\nu}
\]
诱导同构
\[
\widehat{\Omega_{A_\nu|R_\nu}}
\xrightarrow{\ \sim\ }
\widehat{\widehat\Omega_{A_\nu|R_\nu}},
\]
从而
\[
\varprojlim_\nu\widehat{\Omega_{A_\nu|R_\nu}}
\xrightarrow{\ \sim\ }
\varprojlim_\nu
 \widehat{\widehat\Omega_{A_\nu|R_\nu}}.
\]
这里外层帽号表示模关于环所诱导的拓扑的分离完备化. 因此, 若定义
\[
\widehat\Omega_{\mathfrak X_{(\widehat{x})}|S}
=\varprojlim_\nu
 \widehat{\widehat\Omega_{
   \mathcal O_{X_\nu,x}|\mathcal O_{S_\nu,y}}},
\]
则 $\widehat\Omega_{\mathfrak X_{(\widehat{x})}|S}$ 为凝聚 $\mathfrak X_{(\widehat{x})}$-模.

%%%%%%%%%%  scan idx 133  |  folio 122  %%%%%%%%%%

定义 $J_{\mathfrak X_{(\widehat{x})}}^p
 (\mathfrak X_{(\widehat{x})}|S)$ 为 $\mathfrak X_{(\widehat{x})}$ 的闭形式子概形, 由下列理想层定义:
\[
\underline I_{\mathfrak X_{(\widehat{x})}}^p
 (\mathfrak X_{(\widehat{x})}|S)
=\underline I_{\mathfrak X_{(\widehat{x})}}^p
 (\widehat\Omega_{\mathfrak X_{(\widehat{x})}|S}).
\]
于是
\[
J_{\mathfrak X_{(x)}}^p(\mathfrak X_{(x)}|S)_{(\widehat{x})}
\simeq
J_{\mathfrak X_{(\widehat{x})}}^p
 (\mathfrak X_{(\widehat{x})}|S).
\]
若 $x$ 为 $J_{\mathfrak X}^p(\mathfrak X|S)$ 的孤立点, 则自然态射
\[
J_{\mathfrak X_{(\widehat{x})}}^p
 (\mathfrak X_{(\widehat{x})}|S)
\longrightarrow
J_{\mathfrak X_{(x)}}^p(\mathfrak X_{(x)}|S)
\]
为同构.

要构造 $J_{\mathfrak X}^p(\mathfrak X|S)$ 在 $f:\mathfrak X\longrightarrow S$ 下的“像”, 可以直接取下列理想层所定义的 $S$ 的形式子概形:
\[
\ker\left(
 \mathcal O_S\longrightarrow
 \mathcal O_{\mathfrak X}/
 \underline I_{\mathfrak X}^p(\mathfrak X|S)
\right).
\]
但对我们的目的而言, 再次使用 Fitting 理想更为方便.

% END INCLUDED SGA7I_English_current_buildable_source_20260823/components/55_expose_VI_remark_5_8.tex

% BEGIN INCLUDED SGA7I_English_current_buildable_source_20260823/components/56_expose_VI_5_9_through_5_11.tex
\medskip
\noindent\textbf{定义 \SGARefTarget{sga7.sga7i.evi.n5-9}{5.9}.}
设 $S$ 为 诺特形式概形, $f:\mathfrak X\longrightarrow S$ 为本质有限型的 $S$-进形式概形. 假设 $J_{\mathfrak X}^p(\mathfrak X|S)$ 在 $S$ 上紧合, 从而
\[
f_*\mathcal O_{J_{\mathfrak X}^p(\mathfrak X|S)}
\]
为凝聚 $\mathcal O_S$-模. 定义
\[
\underline I_S^p(\mathfrak X|S)
=\underline I_S^0\!\left(
 f_*\mathcal O_{J_{\mathfrak X}^p(\mathfrak X|S)}
\right),
\]
并设 $J_S^p(\mathfrak X|S)$ 为 $\underline I_S^p(\mathfrak X|S)$ 定义的 $S$ 的闭形式子概形.

若存在固定整数 $n$, 对每个 $X_\nu\longrightarrow S_\nu$ 均满足条件 \SGARefLink{sga7.sga7i.evi.n5-4}{5.4}(1) 或 \SGARefLink{sga7.sga7i.evi.n5-4}{5.4}(2), 置
\[
\begin{aligned}
\underline I_{\mathfrak X}(\mathfrak X|S)
 &=\underline I_{\mathfrak X}^n(\mathfrak X|S),
&\qquad
J_{\mathfrak X}(\mathfrak X|S)
 &=J_{\mathfrak X}^n(\mathfrak X|S),\\
\underline I_S(\mathfrak X|S)
 &=\underline I_S^n(\mathfrak X|S),
&
J_S(\mathfrak X|S)
 &=J_S^n(\mathfrak X|S).
\end{aligned}
\]

\medskip
\noindent\textbf{推论 \SGARefTarget{sga7.sga7i.evi.n5-10}{5.10}.}
设 $S$ 为 诺特形式概形, $f:\mathfrak X\longrightarrow S$ 为本质有限型的 $S$-进形式概形, 且 $J_{\mathfrak X}^p(\mathfrak X|S)$ 在 $S$ 上紧合. 则:

\smallskip
\noindent\textbf{(1)}
\[
\bigl|J_S^p(\mathfrak X|S)\bigr|
=f\bigl(\,|J_{\mathfrak X}^p(\mathfrak X|S)|\,\bigr).
\]

%%%%%%%%%%  scan idx 134  |  folio 123  %%%%%%%%%%

\smallskip
\noindent\textbf{(2)} $J_S^p(\mathfrak X|S)$ 的形成与基变换可交换: 若 $\mathfrak X'=\mathfrak X\times_SS'$, 则
\[
J_{S'}^p(\mathfrak X'|S')
=J_S^p(\mathfrak X|S)\times_SS'.
\]

\smallskip
\noindent\textbf{(3)} 若
\[
J_{\mathfrak X}^p(\mathfrak X|S)
=\coprod_{1\leq i\leq m}\mathfrak Y^{(i)}
\]
为不交并, 则
\[
\underline I_S^p(\mathfrak X|S)
=\prod_{1\leq i\leq m}
 \underline I_S^0\bigl(f_*\mathcal O_{\mathfrak Y^{(i)}}\bigr).
\]
特别地, 若
\[
\bigl|J_{X_0}^p(X_0|S_0)\bigr|=\{x_1,\ldots,x_m\}
\]
为离散集, 则
\[
\underline I_S^p(\mathfrak X|S)
=\prod_{1\leq i\leq m}
 \underline I_S^p(\mathfrak X_{(x_i)}|S)
=\prod_{1\leq i\leq m}
 \underline I_S^p(\mathfrak X_{(\widehat{x_i})}|S).
\]

\smallskip
\noindent\textbf{(4)} 设 $S$ 为通常的 诺特概形, $X$ 为本质有限型 $S$-概形, $S_0\hookrightarrow S$ 为闭子概形. 设 $\widehat S$ 为 $S$ 沿 $S_0$ 的形式完备化. 若
\[
\widehat X=X\times_S\widehat S
\]
为 $X$ 沿 $X_0=X\times_SS_0$ 的形式完备化, 则
\[
J_{\widehat S}^p(\widehat X|\widehat S)
=J_S^p(X|S)\times_S\widehat S.
\]
故它是 $J_S^p(X|S)$ 沿 $J_{S_0}^p(X_0|S_0)$ 的形式完备化.

\smallskip
\noindent\textit{证明.}
由 \SGARefLink{sga7.sga7i.evi.n5-1}{5.1}(c),
\[
\operatorname{Supp}
 f_*\mathcal O_{J_{\mathfrak X}^p(\mathfrak X|S)}
=\bigl|J_S^p(\mathfrak X|S)\bigr|,
\]
这就证明了 (1). 断言 (2) 和 (4) 来自 \SGARefLink{sga7.sga7i.evi.n5-1}{5.1}(a), 而 (3) 来自 \SGARefLink{sga7.sga7i.evi.n5-1}{5.1}(d) 以及对孤立点 $x\in J_{\mathfrak X}^p(\mathfrak X|S)$ 成立的同构
\[
J_{\mathfrak X_{(\widehat{x})}}^p
 (\mathfrak X_{(\widehat{x})}|S)
\xrightarrow{\ \sim\ }
J_{\mathfrak X_{(x)}}^p(\mathfrak X_{(x)}|S).
\]
\hfill$\square$

\medskip
\noindent\textbf{注 \SGARefTarget{sga7.sga7i.evi.n5-11}{5.11}.}
我们采用的 Jacobian 理想定义来自底域上的经典定义, 它与基变换可交换, 并在 \SGARefLink{sga7.sga7i.evi.n5-4}{5.4} 的较弱假设下给出光滑性判别. 不过, 在一般情形下, 可能有许多方式定义一个以非光滑轨迹为支集的典范子概形. 下面概述其中一种定义.

对 $A$-代数 $B$, 置
\[
I_{B|A}=\bigcap_E\operatorname{Ann}D^1(B|A,E),
\]
其中 $E$ 遍历所有有限型 $B$-模. 由 \SGARefLink{sga7.sga7i.evi.n3-13}{3.13}(c), 理想 $I_{B|A}$ 的形成与局部化可交换. 因而态射 $f:X\longrightarrow Y$ 决定理想层 $\underline I_{X|Y}$. 若 $f$ 为本质有限型, 则由 \SGARefLink{sga7.sga7i.evi.n3-13}{3.13}(a), $f$ 在 $x$ 处光滑, 当且仅当
\[
x\notin\operatorname{Supp}
 \bigl(\mathcal O_X/\underline I_{X|Y}\bigr).
\]

%%%%%%%%%%  scan idx 135  |  folio 124  %%%%%%%%%%

这一构造自然推广到本质有限型的进性形式概形, 因而也同样适用于形变理论.

% END INCLUDED SGA7I_English_current_buildable_source_20260823/components/56_expose_VI_5_9_through_5_11.tex

% BEGIN INCLUDED SGA7I_English_current_buildable_source_20260823/components/57_expose_VI_section_6_opening_and_lemma_6_1.tex
\par\bigskip
\noindent\textbf{\SGARefAliasTarget{sga7.sga7i.evi.n6-0}6. 非退化二次奇点}

\medskip
本节旨在研究奇点全为非退化二次奇点的 $k$-概形的形变. 由 \SGARefLink{sga7.sga7i.evi.n4-13}{4.13}, 在适当情形下, 这化为一个纯粹局部的问题.

回到 \S4 的记号. 因而 $\Lambda$ 为固定的完备 Noether 局部环, 剩余域为 $k$; $\underline C_\Lambda$ 表示剩余域为 $k$ 的 Artin 局部 $\Lambda$-代数所成的范畴.

设 $X$ 为有限型 $k$-概形, $x$ 为 $X$ 的闭点. 回顾一下, 若
\[
\widehat{\mathcal O}_{X,x}
 \simeq k[[x_1,\ldots,x_n]]/(f)
\]
以及
\[
\left(\frac{\partial f}{\partial x_1},\ldots,
      \frac{\partial f}{\partial x_n}\right)
=(x_1,\ldots,x_n).
\]
则称 $x$ 为\emph{在 $k$ 上有理的非退化二次奇点}. 总可以取 $f$ 为二次型. 事实上:

\medskip
\noindent\textbf{引理 \SGARefTarget{sga7.sga7i.evi.n6-1}{6.1}.}
设 $f\in k[[x_1,\ldots,x_n]]$, 置 $M=(x_1,\ldots,x_n)$. 若
\[
\left(\frac{\partial f}{\partial x_1},\ldots,
      \frac{\partial f}{\partial x_n}\right)=M,
\]
则存在元素 $x_1',\ldots,x_n'$, 满足 $x_i'\equiv x_i\pmod{M^2}$, 使得
\[
f=q(x_1',\ldots,x_n')
\]
其中 $q$ 为二次型.

\smallskip
\noindent\textit{证明.}
使用逐次逼近法. 写成
\[
f=f_2+h_3,
\]
其中 $f_2$ 为二次型, $h_3\in M^3$. 由于 $f$ 的偏导数生成 $M$, 存在 $\varepsilon_1^{(1)},\ldots,\varepsilon_n^{(1)}\in M^2$, 使
\[
\sum_i\varepsilon_i^{(1)}\frac{\partial f}{\partial x_i}
 \equiv h_3\pmod{M^4}.
\]
令 $x_i^{(1)}=x_i-\varepsilon_i^{(1)}$. 则
\[
\begin{aligned}
f(x_1^{(1)},\ldots,x_n^{(1)})
&=f(x_1,\ldots,x_n)
 -\sum_i\varepsilon_i^{(1)}\frac{\partial f}{\partial x_i}
 +\cdots\\
&\equiv f_2(x_1,\ldots,x_n)\pmod{M^4}.
\end{aligned}
\]
故
\[
f(x_1^{(1)},\ldots,x_n^{(1)})
=f_2(x_1,\ldots,x_n)+h_4,
\qquad h_4\in M^4.
\]
此外, 对新变元求得的偏导数仍生成 $M$. 因而可选取 $\varepsilon_1^{(2)},\ldots,\varepsilon_n^{(2)}\in M^3$, 使
\[
\sum_i\varepsilon_i^{(2)}
 \frac{\partial f}{\partial x_i^{(1)}}
 \equiv h_4\pmod{M^5}.
\]

%%%%%%%%%%  scan idx 136  |  folio 125  %%%%%%%%%%

可知
\[
\begin{aligned}
&f(x_1^{(1)}-\varepsilon_1^{(2)},\ldots,
  x_n^{(1)}-\varepsilon_n^{(2)})\\
&\qquad =f(x_1^{(1)},\ldots,x_n^{(1)})
 -\sum_i\varepsilon_i^{(2)}
       \frac{\partial f}{\partial x_i^{(1)}}+\cdots\\
&\qquad \equiv f_2(x_1,\ldots,x_n)\pmod{M^5}.
\end{aligned}
\]
如此继续, 可得元素 $\varepsilon_i^{(\nu)}\in M^{\nu+1}$, 使得置
\[
x_i^{(\nu)}
=x_i-\bigl(\varepsilon_i^{(1)}+\cdots+
            \varepsilon_i^{(\nu)}\bigr),
\]
便有
\[
f(x_1^{(\nu)},\ldots,x_n^{(\nu)})
\equiv f_2(x_1,\ldots,x_n)\pmod{M^{\nu+3}}.
\]
令 $\varepsilon_i=\sum_\nu\varepsilon_i^{(\nu)}$. 由完备性得
\[
f(x_1-\varepsilon_1,\ldots,x_n-\varepsilon_n)
=f_2(x_1,\ldots,x_n),
\]
或等价地,
\[
f(x_1,\ldots,x_n)
=f_2(x_1+\varepsilon_1,\ldots,x_n+\varepsilon_n).
\]
取 $x_i'=x_i+\varepsilon_i$ 即得断言.
\hfill$\square$

% END INCLUDED SGA7I_English_current_buildable_source_20260823/components/57_expose_VI_section_6_opening_and_lemma_6_1.tex

% BEGIN INCLUDED SGA7I_English_current_buildable_source_20260823/components/58_expose_VI_lemma_6_2.tex
% R3 explicit revisions VI-B-C02; see qa/FIDELITY-R3-CHANGES.json.
\par\medskip
\noindent\textbf{引理 \SGARefTarget{sga7.sga7i.evi.n6-2}{6.2}.}
设 $A_0$ 为本质有限型局部 $k$-代数, 满足
\[
\widehat A_0\simeq k[[x_1,\ldots,x_n]]/(f),
\]
其中 $f$ 为非退化二次型. 设 $(R,\mathfrak X)$ 为 $X_0=\operatorname{Spec}(A_0)$ 的形式遍有形变. 则:

\smallskip
\noindent\textbf{(1)} 典范映射
\[
J_{\mathfrak X}(\mathfrak X|R)
 \longrightarrow J_R(\mathfrak X|R)
 \longrightarrow \operatorname{Spec}(\Lambda)
\]
均为同构.

\smallskip
\noindent\textbf{(2)} 理想 $I_R(\mathfrak X|R)$ 为主理想, 且由 $R$ 在 $\Lambda$ 上的一个自由形式生成元生成; 换言之,
\[
I_R(\mathfrak X|R)=(t),
\qquad R=\Lambda[[t]].
\]

\smallskip
\noindent\textbf{(3)} 态射 $\mathfrak X\longrightarrow\operatorname{Spec}(\Lambda)$ 形式光滑.

\smallskip
\noindent\textit{证明.}
写成 $\mathfrak X=\operatorname{Spf}(A)$. 由 \SGARefLink{sga7.sga7i.evi.n4-10}{4.10}(b), $(R,\operatorname{Spf}(\widehat A))$ 是 $\widehat X_0=\operatorname{Spec}(\widehat A_0)$ 的形式遍有形变, 其中
\[
\widehat A=\varprojlim_\nu A\otimes_RR_\nu.
\]
现在
\[
\widehat A_0\simeq k[[x_1,\ldots,x_n]]/(f),
\qquad
\left(\frac{\partial f}{\partial x_1},\ldots,
      \frac{\partial f}{\partial x_n}\right)
=(x_1,\ldots,x_n).
\]
立即得出
\[
\bigl[\underline M_{\widehat X_0}(k[\varepsilon])\bigr]
\simeq k.
\]
从而, 由 \SGARefLink{sga7.sga7i.evi.n4-14}{4.14},
\[
R\simeq\Lambda[[t]],
\qquad
\widehat A\simeq R[[x_1,\ldots,x_n]]/(F-t),
\]
其中 $F$ 是将 $f$ 的系数提升到 $\Lambda$ 而得到的 $\Lambda[x_1,\ldots,x_n]\subset R[[x_1,\ldots,x_n]]$ 中的二次多项式.

序列
\[
0\longrightarrow\widehat A
 \xrightarrow{\ dF\ }
 \widehat A\,dx_1\oplus\cdots\oplus\widehat A\,dx_n
 \longrightarrow\widehat\Omega_{\widehat A|R}\longrightarrow0
\]
为正合列, 其中
\[
1\longmapsto
 \frac{\partial F}{\partial x_1}dx_1+\cdots+
 \frac{\partial F}{\partial x_n}dx_n.
\]

%%%%%%%%%%  scan idx 137  |  folio 126  %%%%%%%%%%

因此
\[
I_{\widehat A}(\widehat A|R)
=\left(\frac{\partial F}{\partial x_1},\ldots,
       \frac{\partial F}{\partial x_n}\right)
=(x_1,\ldots,x_n),
\]
从而
\[
\begin{aligned}
\widehat A/I_{\widehat A}(\widehat A|R)
&=\widehat A/(x_1,\ldots,x_n)\\
&=R[[x_1,\ldots,x_n]]/(x_1,\ldots,x_n,F-t)
 \simeq\Lambda.
\end{aligned}
\]
故典范映射 $\Lambda\longrightarrow
\widehat A/I_{\widehat A}(\widehat A|R)$ 为同构. 完备化比较将此商等同于 $A/I_A(A|R)$, 因而
\[
\Lambda\xrightarrow{\ \sim\ }R/I_R(A|R).
\]
这就证明了 (1) 和 (2). 最后, (3) 来自如下 $\Lambda$-代数同构:
\[
\begin{aligned}
\widehat A
&\simeq R[[x_1,\ldots,x_n]]/(F-t)\\
&\simeq\Lambda[[t,x_1,\ldots,x_n]]/(F-t)
 \simeq\Lambda[[x_1,\ldots,x_n]]
\end{aligned}
\]

\hfill$\square$

% END INCLUDED SGA7I_English_current_buildable_source_20260823/components/58_expose_VI_lemma_6_2.tex

% BEGIN INCLUDED SGA7I_English_current_buildable_source_20260823/components/59_expose_VI_corollary_6_3.tex
\par\medskip
\noindent\textbf{推论 \SGARefTarget{sga7.sga7i.evi.n6-3}{6.3}.}\SGAEditorialMark{editorial-vi-square-root}{校记 B}
设 $X_0$ 为有限型 $k$-概形, $x\in X_0$ 为在 $k$ 上有理的非退化二次奇点.

\smallskip
\noindent\textbf{(1)} 设 $(R,\mathfrak X)$ 与 $(R,\mathfrak Y)$ 是 $X_0$ 的两个形变, 其中 $R\in\operatorname{ob}(\widehat{\underline C}_\Lambda)$. 假设 $k=k^2$. 则
\[
\mathfrak X_{(\widehat x)}\simeq
\mathfrak Y_{(\widehat x)}
\]
作为 $R$-进形式概形成立, 当且仅当
\[
I_R(\mathfrak X_{(x)}|R)=I_R(\mathfrak Y_{(x)}|R).
\]

\smallskip
\noindent\textbf{(2)} 设 $R\in\operatorname{ob}(\widehat{\underline C}_\Lambda)$ 为正则局部环, $(R,\mathfrak X)$ 为 $X_0$ 的形变. 则 $\mathcal O_{\mathfrak X,x}$ 正则, 当且仅当
\[
I_R(\mathfrak X_{(x)}|R)\not\subset\mathfrak m_R^2;
\]
等价地, $I_R(\mathfrak X_{(x)}|R)$ 是 $\operatorname{Spf}(R)$ 中的正则除子.

\smallskip
\noindent\textit{证明.}
设 $(R_0,Z)$ 为 $\operatorname{Spec}(\widehat{\mathcal O}_{X_0,x})$ 的形式遍有形变. 由于 $x$ 为在 $k$ 上有理的非退化二次奇点, 由引理 \SGARefLink{sga7.sga7i.evi.n6-2}{6.2} 可取
\[
R_0\simeq\Lambda[[t]],
\qquad
Z\simeq\operatorname{Spf}
 \bigl(R_0[[X_1,\ldots,X_n]]/(F-t)\bigr).
\]

对于 (1), 形变 $(R,\mathfrak X_{(\widehat x)})$ 和 $(R,\mathfrak Y_{(\widehat x)})$ 由 $(R_0,Z)$ 通过连续局部 $\Lambda$-代数同态
\[
\varphi,\psi:R_0\longrightarrow R.
\]
诱导. 令 $a=\varphi(t)$, $b=\psi(t)$. 则
\[
\begin{aligned}
\mathfrak X_{(\widehat x)}
&\simeq\operatorname{Spf}
 \bigl(R[[X_1,\ldots,X_n]]/(F-a)\bigr),\\
\mathfrak Y_{(\widehat x)}
&\simeq\operatorname{Spf}
 \bigl(R[[X_1,\ldots,X_n]]/(F-b)\bigr).
\end{aligned}
\]

%%%%%%%%%%  scan idx 138  |  folio 127  %%%%%%%%%%

假设 $I_R(\mathfrak X_{(x)}|R)=I_R(\mathfrak Y_{(x)}|R)$; 等价地, $R$ 中 $(a)=(b)$. 则存在单位 $u\in R$ 使 $a=ub$. 由于 $k=k^2$ 且 $R$ 完备, 存在单位 $v\in R$ 使 $u=v^2$, 从而 $a=v^2b$. 于是
\[
(F-a)\longmapsto v^2(F-b)
\]
在 $X_i\mapsto vX_i$ 下成立; 这一代换诱导 $R$-代数同构
\[
R[[X_1,\ldots,X_n]]/(F-a)
 \xrightarrow{\ \sim\ }
R[[X_1,\ldots,X_n]]/(F-b),
\]
因为 $F$ 为二次型. 反向结论显然.

对于 (2), 存在连续局部 $\Lambda$-代数同态
\[
\varphi:R_0=\Lambda[[t]]\longrightarrow R
\]
使得
\[
\mathfrak X_{(\widehat x)}
\simeq\operatorname{Spf}
 \bigl(R[[X_1,\ldots,X_n]]/(F-\varphi(t))\bigr).
\]
此外,
\[
I_R(\mathfrak X_{(x)}|R)
=I_R(\mathfrak X_{(\widehat x)}|R)
=(\varphi(t)),
\]
以及
\[
R[[X_1,\ldots,X_n]]/(F-\varphi(t))
=\widehat{\mathcal O}_{\mathfrak X,x}
\]
正则, 当且仅当 $\mathcal O_{\mathfrak X,x}$ 正则. 由于 $R$ 正则, $R[[X_1,\ldots,X_n]]$ 为正则局部环. 若以 $M$ 表示其极大理想, 则它模 $F-\varphi(t)$ 的商正则, 当且仅当
\[
F-\varphi(t)\notin M^2,
\]
这等价于 $\varphi(t)\notin\mathfrak m_R^2$, 从而等价于
\[
I_R(\mathfrak X_{(x)}|R)\not\subset\mathfrak m_R^2.
\]
\hfill$\square$

% END INCLUDED SGA7I_English_current_buildable_source_20260823/components/59_expose_VI_corollary_6_3.tex

% BEGIN INCLUDED SGA7I_English_current_buildable_source_20260823/components/60_expose_VI_theorem_6_4_through_part_4.tex
% R3 explicit revisions VI-B-C07; see qa/FIDELITY-R3-CHANGES.json.
\par\medskip
\noindent\textbf{定理 \SGARefTarget{sga7.sga7i.evi.n6-4}{6.4}.}
设 $X_0$ 为有限型 $k$-概形, 其非光滑点孤立, $(R,\mathfrak X)$ 为 $X_0$ 的形式遍有形变. 记 $X_0$ 的非光滑点集合为 $\{z_1,\ldots,z_m\}$. 假设
\[
H^2(X_0,\underline D^0_{X_0})=0
\]
且每个 $z_i$ 均为在 $k$ 上有理的非退化二次奇点. 则:

\smallskip
\noindent\textbf{(1)} $R$ 在 $\Lambda$ 上形式光滑; 即 $R$ 是 $\Lambda$ 上的形式幂级数环.

\smallskip
\noindent\textbf{(2)} 有不交并分解
\[
J_{\mathfrak X}(\mathfrak X|R)
=\coprod_{1\leq i\leq m}
 J_{\mathfrak X_{(z_i)}}(\mathfrak X_{(z_i)}|R),
\]
且典范态射
\[
J_{\mathfrak X_{(z_i)}}(\mathfrak X_{(z_i)}|R)
 \xrightarrow{\ \sim\ }
J_R(\mathfrak X_{(z_i)}|R)
\]
为同构.

\smallskip
\noindent\textbf{(3)} 各理想 $I_R(\mathfrak X_{(z_i)}|R)$ 均为主理想, 下列理想也是主理想:
\[
I_R(\mathfrak X|R)
=\prod_{1\leq i\leq m}I_R(\mathfrak X_{(z_i)}|R).
\]
若 $I_R(\mathfrak X_{(z_i)}|R)=(t_i)$, 则 $t_1,\ldots,t_m$ 可以扩充为 $R$ 在 $\Lambda$ 上的一组自由形式生成元.

\smallskip
\noindent\textbf{(4)} 态射 $\mathfrak X\longrightarrow\operatorname{Spec}(\Lambda)$ 形式光滑; 等价地,
\[
J_{\mathfrak X}(\mathfrak X|\Lambda)=\varnothing.
\]

%%%%%%%%%%  scan idx 139  |  folio 128  %%%%%%%%%%

\smallskip
\noindent\textbf{(5)} 若 $X_0$ 为完备不可约曲线, 则
\[
\dim\underline M_{X_0}=3g-3+a,
\]
其中
\[
g=\dim_kH^1(X_0,\mathcal O_{X_0}),
\qquad
a=\dim_kH^0(X_0,\underline D^0_{X_0}).
\]

\smallskip
\noindent\textit{证明.}
断言 (1) 来自 (2) 和 (3). (2) 中的第一个断言来自 \SGARefLink{sga7.sga7i.evi.n5-7}{5.7}(2). 对每个 $i$, 设 $(R_i,\mathfrak Y_i)$ 为 $\operatorname{Spec}(\widehat{\mathcal O}_{X_0,z_i})$ 的形式遍有形变. 由于 $(R,\mathfrak X_{(\widehat z_i)})$ 是该局部概形的一个形变, 存在连续局部 $\Lambda$-代数同态
\[
\varphi_i:R_i\longrightarrow R
\]
及笛卡尔图
\[
\begin{tikzcd}[column sep=huge,row sep=large]
\mathfrak X_{(\widehat z_i)}
  \arrow[r] \arrow[d]
& \mathfrak Y_i \arrow[d]\\
\operatorname{Spf}(R)
  \arrow[r,"{\operatorname{Spf}(\varphi_i)}"']
& \operatorname{Spf}(R_i).
\end{tikzcd}
\]
因此,
\[
\begin{aligned}
J_{\mathfrak X_{(z_i)}}(\mathfrak X_{(z_i)}|R)
&=J_{\mathfrak X_{(\widehat z_i)}}
   (\mathfrak X_{(\widehat z_i)}|R)\\
&\simeq
J_{\mathfrak Y_i}(\mathfrak Y_i|R_i)
 \times_{\operatorname{Spf}(R_i)}\operatorname{Spf}(R),
\end{aligned}
\]
以及
\[
\begin{aligned}
J_R(\mathfrak X_{(z_i)}|R)
&=J_R(\mathfrak X_{(\widehat z_i)}|R)\\
&\simeq
J_{R_i}(\mathfrak Y_i|R_i)
 \times_{\operatorname{Spf}(R_i)}\operatorname{Spf}(R).
\end{aligned}
\]
由引理 \SGARefLink{sga7.sga7i.evi.n6-2}{6.2}(1),
\[
J_{\mathfrak Y_i}(\mathfrak Y_i|R_i)
 \longrightarrow J_{R_i}(\mathfrak Y_i|R_i)
\]
为同构, 这就证明了 (2) 中余下的断言.

由引理 \SGARefLink{sga7.sga7i.evi.n6-2}{6.2}(2),
\[
I_{R_i}(\mathfrak Y_i|R_i)=(u_i),
\qquad R_i=\Lambda[[u_i]].
\]
可知
\[
I_R(\mathfrak X_{(z_i)}|R)
=I_R(\mathfrak X_{(\widehat z_i)}|R)
=(t_i),
\qquad t_i=\varphi_i(u_i),
\]
因此
\[
I_R(\mathfrak X|R)
=\prod_{1\leq i\leq m}I_R(\mathfrak X_{(z_i)}|R)
=\left(\prod_{i=1}^m t_i\right).
\]
由 \SGARefLink{sga7.sga7i.evi.n4-13}{4.13}, 自然态射
\[
\underline M_{X_0}\longrightarrow
\underline M_{X_0,(\widehat z_1)}\times\cdots\times
\underline M_{X_0,(\widehat z_m)}
\]
光滑. 因而
\[
R_1\widehat\otimes_\Lambda\cdots
 \widehat\otimes_\Lambda R_m\longrightarrow R
\]
形式光滑. 故 $t_1,\ldots,t_m$ 可扩充为 $R$ 在 $\Lambda$ 上的一组自由形式生成元, 这就证明了 (1) 和 (3).

%%%%%%%%%%  scan idx 140  |  folio 129  %%%%%%%%%%

对于 (4), 投影
\[
\mathfrak X_{(\widehat z_i)}
 \xrightarrow{\ \sim\ }
\mathfrak Y_i\times_{\operatorname{Spf}(R_i)}
 \operatorname{Spf}(R)
 \longrightarrow\mathfrak Y_i
\]
形式光滑, 因为 $\operatorname{Spf}(R)\longrightarrow\operatorname{Spf}(R_i)$ 形式光滑. 引理 \SGARefLink{sga7.sga7i.evi.n6-2}{6.2}(3) 表明 $\mathfrak Y_i\longrightarrow\operatorname{Spec}(\Lambda)$ 形式光滑. 因而每个 $\mathfrak X_{(\widehat z_i)}\longrightarrow
\operatorname{Spec}(\Lambda)$ 都形式光滑. 由于各 $z_i$ 恰为 $X_0$ 的全部非光滑点, 可知 $\mathfrak X\longrightarrow\operatorname{Spec}(\Lambda)$ 形式光滑.

% END INCLUDED SGA7I_English_current_buildable_source_20260823/components/60_expose_VI_theorem_6_4_through_part_4.tex

% BEGIN INCLUDED SGA7I_English_current_buildable_source_20260823/components/61_expose_VI_theorem_6_4_part_5_and_remark_6_5.tex
% R3 explicit revisions VI-B-C01; see qa/FIDELITY-R3-CHANGES.json.
对于 (5), $R$ 在 $\Lambda$ 上形式光滑, 且由 \SGARefLink{sga7.sga7i.evi.n4-4}{4.4} 和 \SGARefLink{sga7.sga7i.evi.n4-11}{4.11}, 序列
\[
0\longrightarrow H^1(X_0,\underline D^0_{X_0})
 \longrightarrow[\underline M_{X_0}(k[\varepsilon])]
 \longrightarrow H^0(X_0,\underline D^1_{X_0})
 \longrightarrow0
\]
为正合列. 因此
\[
\begin{aligned}
\dim\underline M_{X_0}
&=\dim_kH^1(X_0,\underline D^0_{X_0})
 +\dim_kH^0(X_0,\underline D^1_{X_0})\\
&=\dim_kH^1(X_0,\underline D^0_{X_0})+m.
\end{aligned}
\]
因而只需证明
\[
\dim_kH^1(X_0,\underline D^0_{X_0})+m
=3g-3+\dim_kH^0(X_0,\underline D^0_{X_0}),
\]
或等价地
\[
\chi(\underline D^0_{X_0})=3-3g+m,
\]
其中对凝聚 $\mathcal O_{X_0}$-模 $\mathcal E$, 置
\[
\chi(\mathcal E)
=\dim_kH^0(X_0,\mathcal E)-\dim_kH^1(X_0,\mathcal E).
\]

将底域扩张为其代数闭包不会改变 $\underline D^0_{X_0}$ 的 Euler 示性数, 因而可假设 $k$ 为代数闭域. 设 $z$ 为 $X_0$ 的一个非光滑点, 置
\[
A=\widehat{\mathcal O}_{X_0,z},
\]
并设 $A'$ 为 $A$ 的正规化. 按假设,
\[
A\simeq k[[x,y]]/(f),
\]
其中 $f$ 为非退化二次型. 在代数闭域 $k$ 上可取 $f=xy$, 从而
\[
A\simeq k[[x,y]]/(xy).
\]
立即得出
\[
\ell(A'/A)=1,
\qquad
\operatorname{Der}_k(A,\mathfrak m_A)
=\operatorname{Der}_k(A,A).
\]
设 $C$ 为 $A$ 在 $A'$ 中的导理想 (conductor). 由于 $A$ 为完全交,
\[
\ell(A/C)=\ell(A'/A)=1,
\]
从而 $C=\mathfrak m_A$. 因此
\[
A'=\{a'\in K\mid a'\mathfrak m_A\subset\mathfrak m_A\},
\]
其中 $K$ 为 $A$ 的全分式环.

断言
\[
\operatorname{Der}_k(A,A)
 \subset\operatorname{Der}_k(A',A')
\]
在 $\operatorname{Der}_k(K,K)$ 中成立. 事实上, 设 $D\in\operatorname{Der}_k(A,A)$. 对 $a'\in A'$ 和 $x\in\mathfrak m_A$, 有 $xa'\in\mathfrak m_A$, 故 $D(xa')\in\mathfrak m_A$, 且 $a'D(x)\in\mathfrak m_A$. 因而等式
\[
D(xa')=a'D(x)+xD(a')
\]
给出 $xD(a')\in\mathfrak m_A$, 对每个 $x\in\mathfrak m_A$ 都成立, 从而 $D(a')\in A'$.

%%%%%%%%%%  scan idx 141  |  folio 130  %%%%%%%%%%

此外, 正合列
\[
0\longrightarrow\operatorname{Der}_k(A,A)
 \longrightarrow\operatorname{Der}_k(A',A')
 \longrightarrow A'/\mathfrak m_AA'\longrightarrow0
\]
可写成
\[
0\longrightarrow\operatorname{Der}_k(A,A)
 \longrightarrow\operatorname{Der}_k(A',A')
 \longrightarrow A'/C\longrightarrow0.
\]
若 $X_0'$ 为 $X_0$ 的正规化, 则这些局部正合列给出正合列
\[
0\longrightarrow\underline D^0_{X_0}
 \longrightarrow\underline D^0_{X_0'}
 \longrightarrow\mathcal O_{X_0'}/C\longrightarrow0.
\]
由于 $X_0$ 局部为完全交,
\[
\ell(\mathcal O_{X_0'}/C)
=2\ell(\mathcal O_{X_0'}/\mathcal O_{X_0})=2m,
\]
而 Riemann--Roch 定理给出
\[
\chi(\underline D^0_{X_0'})=3-3(g-m).
\]
因此,
\[
\begin{aligned}
\chi(\underline D^0_{X_0})
&=\chi(\underline D^0_{X_0'})
 -\ell(\mathcal O_{X_0'}/C)\\
&=3-3(g-m)-2m=3-3g+m.
\end{aligned}
\]
证明完毕.
\hfill$\square$

\par\medskip
\noindent\textbf{注 \SGARefTarget{sga7.sga7i.evi.n6-5}{6.5}.}
在定理 \SGARefLink{sga7.sga7i.evi.n6-4}{6.4}(5) 的公式中, 数 $a=\dim_kH^0(X_0,\underline D^0_{X_0})$ 容易由 Riemann--Roch 定理计算:
\[
a=
\begin{cases}
0,&g\geq2,\\
1,&g=1,\\
3,&g=0.
\end{cases}
\]

最后说明一下特征 $2$ 的情形. 若 $\operatorname{char}k=2$, 则奇数个变元的二次型不可能非退化. 等价地, 特征 $2$ 时不存在偶数维的非退化二次奇点. 此时可改为考虑退化程度最小的二次型.

% END INCLUDED SGA7I_English_current_buildable_source_20260823/components/61_expose_VI_theorem_6_4_part_5_and_remark_6_5.tex

% BEGIN INCLUDED SGA7I_English_current_buildable_source_20260823/components/62_expose_VI_6_6_through_6_8.tex
%%%%%%%%%%  scan idx 142  |  folio 131  %%%%%%%%%%

\par\medskip
\noindent\textbf{定义 \SGARefTarget{sga7.sga7i.evi.n6-6}{6.6}.}
设 $X$ 为有限型 $k$-概形. 若满足下列条件, 则称闭点 $z\in X$ 为\emph{在 $k$ 上有理的普通二次奇点}.

\smallskip
\noindent\textbf{(i)} 若 $\operatorname{char}k\neq2$, 或 $\operatorname{char}k=2$ 且 $\dim\mathcal O_{X,z}$ 为奇数, 则 $z$ 为在 $k$ 上有理的非退化二次奇点.

\smallskip
\noindent\textbf{(ii)} 若 $\operatorname{char}k=2$ 且 $\dim\mathcal O_{X,z}$ 为偶数, 则
\[
\overline k\otimes_k\widehat{\mathcal O}_{X,z}
 \simeq
\overline k[[X_0,X_1,\ldots,X_{2n}]]/(f),
\]
其中 $\overline k$ 为 $k$ 的代数闭包, 且
\[
f=X_0^2+X_1X_2+X_3X_4+\cdots+X_{2n-1}X_{2n}.
\]

\par\medskip
\noindent\textbf{引理 \SGARefTarget{sga7.sga7i.evi.n6-7}{6.7}.}\SGAEditorialMark{editorial-vi-mixedchar}{校记 A}
设 $k$ 为特征 $2$ 的域, 置
\[
A_0=k[[X_0,X_1,\ldots,X_{2n}]]/(f),
\qquad
f=X_0^2+X_1X_2+X_3X_4+\cdots+X_{2n-1}X_{2n}.
\]
设 $(R,\mathfrak X)$ 为 $X_0=\operatorname{Spec}(A_0)$ 的形式遍有形变. 则
\[
R=\Lambda[[t,t_1]],
\]
其中 $t,t_1$ 是 $R$ 在 $\Lambda$ 上的自由形式生成元, 且
\[
I_R(\mathfrak X|R)=(t^2).
\]
此外, $\mathfrak X\longrightarrow\operatorname{Spec}(\Lambda)$ 形式光滑.

\smallskip
\noindent\textit{证明.}
写成 $\mathfrak X=\operatorname{Spf}(A)$. 由于
\[
\left(f,\frac{\partial f}{\partial X_0},\ldots,
         \frac{\partial f}{\partial X_{2n}}\right)
=(X_0^2,X_1,\ldots,X_{2n}),
\]
有
\[
D^1(A_0|k,A_0)\simeq k[X_0]/(X_0^2).
\]
故 $[\underline M_{X_0}(k[\varepsilon])]$ 为二维 $k$-向量空间. 由 \SGARefLink{sga7.sga7i.evi.n4-14}{4.14},
\[
R\simeq\Lambda[[t,t_1]],
\qquad
A\simeq R[[X_0,X_1,\ldots,X_{2n}]]/(F),
\]
其中
\[
F=f+t_1+tX_0.
\]
正合列
\[
0\longrightarrow A\xrightarrow{\ dF\ }
A\,dX_0\oplus\cdots\oplus A\,dX_{2n}
\longrightarrow\widehat\Omega_{A|R}\longrightarrow0,
\]
其中
\[
dF:\ 1\longmapsto
\sum_{0\leq i\leq2n}\frac{\partial f}{\partial X_i}\,dX_i
+t\,dX_0,
\]
给出
\[
I_A(A|R)
=(t,X_1,\ldots,X_{2n})
=(X_0^2+t_1,t,X_1,\ldots,X_{2n}).
\]
% The printed formula omits the coefficient t before dX_0, although
% F=f+t_1+tX_0 and the ensuing ideal both force this derivative term.

%%%%%%%%%%  scan idx 143  |  folio 132  %%%%%%%%%%

从而
\[
\begin{aligned}
A/I_A(A|R)
&\simeq A/(X_0^2+t_1,t,X_1,\ldots,X_{2n})\\
&\simeq\Lambda[[t_1,X_0]]/(X_0^2+t_1).
\end{aligned}
\]
作为 $R$-模, 最后这个模同构于 $R/(t)\oplus R/(t)$. 可知
\[
I_R(A|R)=(t^2),
\]
这就证明了第一个断言. 另一方面,
\[
\begin{aligned}
A
&\simeq
\Lambda[[t,t_1,X_0,X_1,\ldots,X_{2n}]]/
 (X_0^2+X_1X_2+\cdots+X_{2n-1}X_{2n}+tX_0+t_1)\\
&\simeq\Lambda[[t,X_0,X_1,\ldots,X_{2n}]],
\end{aligned}
\]
故 $A$ 在 $\Lambda$ 上形式光滑.
\hfill$\square$

引理 \SGARefLink{sga7.sga7i.evi.n6-7}{6.7} 给出下列与定理 \SGARefLink{sga7.sga7i.evi.n6-4}{6.4} 类似的结果; 证明留给读者.

\par\medskip
\noindent\textbf{推论 \SGARefTarget{sga7.sga7i.evi.n6-8}{6.8}.}\SGAEditorialMark{editorial-vi-mixedchar}{校记 A}
设 $X_0$ 为有限型 $k$-概形, 其非光滑点孤立, 假设
\[
H^2(X_0,\underline D^0_{X_0})=0,
\]
并设 $(R,\mathfrak X)$ 为 $X_0$ 的形式遍有形变. 记非光滑点集合为 $\{z_1,\ldots,z_m\}$. 假设 $\operatorname{char}k=2$, $\dim X_0$ 为偶数, 且各 $z_i$ 均为普通二次奇点. 则:

\smallskip
\noindent\textbf{(1)} $R$ 在 $\Lambda$ 上形式光滑.

\smallskip
\noindent\textbf{(2)} 各理想 $I_R(\mathfrak X_{(z_i)}|R)$ 均为主理想, 下列理想也是主理想:
\[
I_R(\mathfrak X|R)
=\prod_{1\leq i\leq m}I_R(\mathfrak X_{(z_i)}|R).
\]
更确切地,
\[
I_R(\mathfrak X_{(z_i)}|R)=(t_i^2),
\]
且 $t_1,\ldots,t_m$ 可以扩充为 $R$ 在 $\Lambda$ 上的一组自由形式生成元.

\smallskip
\noindent\textbf{(3)} 态射 $\mathfrak X\longrightarrow\operatorname{Spec}(\Lambda)$ 形式光滑.

\smallskip
\noindent\textbf{(4)} 若 $k$ 为代数闭域, $X_0$ 为完备不可约曲线, 则
\[
\dim\underline M_{X_0}=3g-3+a,
\]
其中
\[
g=\dim_kH^1(X_0,\mathcal O_{X_0}),
\qquad
a=\dim_kH^0(X_0,\underline D^0_{X_0}).
\]

% END INCLUDED SGA7I_English_current_buildable_source_20260823/components/62_expose_VI_6_6_through_6_8.tex
